\documentclass[12pt]{article}
\usepackage{amsmath}
\usepackage{amsfonts}
\usepackage{float}
\usepackage{fancyhdr}
\usepackage{graphicx}
\usepackage{setspace}
\usepackage[colorlinks=true,linkcolor=blue, citecolor=red]{hyperref}
\usepackage{url}
\usepackage[top=.75in, left=.75in, right=.75in, bottom=1in]{geometry}

% For algorithm
\usepackage{algorithm}
\usepackage{algpseudocode}

% ============ CODE ============
\usepackage{listings}
\usepackage{xcolor}
\definecolor{codegreen}{rgb}{0,0.6,0}
\definecolor{codegray}{rgb}{0.5,0.5,0.5}
\definecolor{codepurple}{rgb}{0.58,0,0.82}
\definecolor{backcolour}{rgb}{0.95,0.95,0.92}

% Styling for the code.
\lstdefinestyle{mystyle}{
    backgroundcolor=\color{backcolour},   
    commentstyle=\color{codegreen},
    keywordstyle=\color{magenta},
    numberstyle=\tiny\color{codegray},
    stringstyle=\color{codepurple},
    basicstyle=\ttfamily\footnotesize,
    breakatwhitespace=false,         
    breaklines=true,                 
    captionpos=b,                    
    keepspaces=true,                 
    numbers=left,                    
    numbersep=5pt,                  
    showspaces=false,                
    showstringspaces=false,
    showtabs=false,                  
    tabsize=2
}
\lstset{style=mystyle}

% Table packages and Unicode / Font support
\usepackage{array}
\usepackage{longtable}
\usepackage{enumitem}
\usepackage{booktabs}
\usepackage[utf8]{inputenc}
\DeclareUnicodeCharacter{2264}{\ensuremath{\le}}
\DeclareUnicodeCharacter{2194}{\ensuremath{\leftrightarrow}}
\usepackage[T1]{fontenc}
\usepackage{lmodern}
\graphicspath{{../../../../templates/}{../../prototype/}{figures/}{./}}
\setlength{\emergencystretch}{3em}
\setlist[enumerate]{label=\arabic*.,leftmargin=*}
\setlist[itemize]{leftmargin=*}

% ============ HEADER AND FOOTER ============
\setlength{\headheight}{29.43912pt}

\pagestyle{fancy}
\fancyfoot{}
\fancyfoot[R]{Page \thepage}

\setlength{\parindent}{0pt}
\renewcommand{\baselinestretch}{1.5}

\newcommand{\coursename}{Introduction of Software Engineering}
\newcommand{\reportname}{Studitopia}
\newcommand{\reporttitle}{\href{https://github.com/RoundFloor/Studitopia.git}{Studitopia}}

\newcommand{\studentname}{%
Nguyen Thanh An (24125024)\\
Nguyen Chi Bao (24125025)\\
Tran Nguyen Khanh Binh (24125026)\\
Tran Canh Anh Tuan (24125048)\\
Nguyen Le Quoc Huy (24125100)
}
\newcommand{\teachername}{Ho Tuan Thanh, Nguyen Thi Minh Tuyen}

% Header
\lhead{\reporttitle}
\rhead{
University of Science - VNUHCM\\
\coursename
}

% Footer
\newcommand{\leftfooter}{}
\lfoot{\leftfooter}

% ============ DOCUMENT ============
\begin{document}

\pagenumbering{roman}
\begin{titlepage}
\newcommand{\HRule}{\rule{\linewidth}{0.5mm}}
\centering

\textsc{\LARGE VNU-HCM}\\[0.5cm]
\textsc{\Large University of Science}\\[0.5cm]
\textsc{\large Faculty of Information Technology}\\[0.5cm]
\textsc{Computer Science}\\[0.5cm]

\HRule \\[0.4cm]
{ 
\huge{\bfseries{\reporttitle}}\\[0.5cm]
\large{\bfseries{Forum Interaction --- Use Case Diagram}}
}\\[0.4cm]
\HRule \\[0.5cm]

\textbf{\large Subject: \coursename}\\[0.5cm]

\begin{minipage}[t]{0.55\textwidth}
\begin{flushleft} \large
\emph{Students:}\\
\studentname
\end{flushleft}
\end{minipage}
~
\begin{minipage}[t]{0.4\textwidth}
\begin{flushright} \large
\emph{Supervising Teachers:} \\
\teachername
\end{flushright}
\end{minipage}\\[1cm]

{\large \today}\\[1cm]

\includegraphics[scale=.20]{img/hcmus-logo.png}\\[1cm] 

\vfill
\end{titlepage}	

\tableofcontents
\pagebreak

\pagenumbering{arabic}
\setcounter{page}{1}

\hypertarget{sec:forum-interaction-use-case-diagram}{}
\addcontentsline{toc}{section}{Forum Interaction --- Use Case Diagram}
\section*{Forum Interaction --- Use Case Diagram}
\par\medskip

\textbf{Version:} \texttt{1.0}
\par\medskip

\textbf{Performed by:} [Nguyen Le Quoc Huy] | \textbf{Edited by:} [Nguyen Le Quoc Huy]
\par\medskip

Merged from the focused use case diagrams in this functional group.
\par\medskip

\begin{center}
\includegraphics[width=0.85\textwidth,height=0.75\textheight,keepaspectratio]{"figures/diagram_1.png"}
\end{center}
\par\medskip

\textbf{Diagram notation.} Person shape = human actor; cloud = supporting system; blue = base use case; purple = included; green = extending; solid line = association; dashed arrow = include or extend; labelled solid arrow = generalization.
\par\medskip

\pagebreak
\hypertarget{sec:uc-05-comment-and-answer}{}
\addcontentsline{toc}{section}{UC-05 --- Comment and Answer}
\section*{UC-05 --- Comment and Answer}
\par\medskip

\textbf{Version:} \texttt{2.0}
\par\medskip

\textbf{Performed by:} [Tran Nguyen Khanh Binh] | \textbf{Edited by:} [Nguyen Le Quoc Huy]
\par\medskip

\begingroup
\small
\renewcommand{\baselinestretch}{1.0}\selectfont
\renewcommand{\arraystretch}{1.15}
\setlength{\tabcolsep}{6pt}
\begin{longtable}{|>{\raggedright\arraybackslash}p{0.09\textwidth}|>{\raggedright\arraybackslash}p{0.84\textwidth}|}
\hline
\textbf{Field} & \textbf{Value} \\ \hline
\endhead
Use Case ID & UC-05 \\ \hline
Name & Comment and Answer \\ \hline
Primary Actor & University Student \\ \hline
Supporting Actors & Post Author; Notification Service \\ \hline
Feature Group & Forum (Primary) \\ \hline
Priority & High \\ \hline
Related User Stories & US-18 \\ \hline
Related Requirements & US-18, US-17, US-20, US-38; FR-18.1 to FR-18.6, FR-17.3, FR-20.1 to FR-20.5, FR-38.1; NFR-PERF-4, NFR-PERF-6, NFR-SEC-9 \\ \hline
Source Documents & \texttt{proposals/use\_cases\_model/UC-05-comment-answer.md}; \texttt{proposals/use\_cases\_description.md} (UC-05); \texttt{proposals/use\_cases\_model.md}; \texttt{proposals/use\_case\_diagrams.md}; \texttt{PA3-Group05/use\_case.md} (section 5); \texttt{proposals/overview.md}, \texttt{functional\_requirement.md}, \texttt{non\_functional\_requirement.md}, \texttt{userstory.md}, \texttt{feature\_flow.md}, \texttt{screen\_flow.md}, \texttt{API.md}, \texttt{response\_code.md}, \texttt{post\_management.md}, \texttt{moderation\_queue.md}, \texttt{list\_post.md}, \texttt{settings.md}, \texttt{friend.md}, \texttt{friend\_profile.md} \\ \hline
\end{longtable}
\endgroup
\par\medskip

\hypertarget{sec:part-a-use-case-model}{}
\addcontentsline{toc}{subsection}{Part A --- Use Case Model}
\subsection*{Part A --- Use Case Model}
\par\medskip

\hypertarget{sec:a1-focused-use-case-diagram}{}
\addcontentsline{toc}{subsubsection}{A1. Focused Use Case Diagram}
\subsubsection*{A1. Focused Use Case Diagram}
\par\medskip

\begin{center}
\includegraphics[width=0.85\textwidth,height=0.75\textheight,keepaspectratio]{"figures/UC-05-comment-answer_1.png"}
\end{center}
\par\medskip

\textbf{Diagram notation.} Person shape = human actor; cloud = supporting system; blue = selected base use case; purple = included; green = extending; solid line = association; dashed arrow = include or extend; labelled solid arrow = generalization.
\par\medskip

\hypertarget{sec:a2-actor-register}{}
\addcontentsline{toc}{subsubsection}{A2. Actor Register}
\subsubsection*{A2. Actor Register}
\par\medskip

\begingroup
\small
\renewcommand{\baselinestretch}{1.0}\selectfont
\renewcommand{\arraystretch}{1.15}
\setlength{\tabcolsep}{6pt}
\begin{longtable}{|>{\raggedright\arraybackslash}p{0.10\textwidth}|>{\raggedright\arraybackslash}p{0.10\textwidth}|>{\raggedright\arraybackslash}p{0.72\textwidth}|}
\hline
\textbf{Actor} & \textbf{Type} & \textbf{Description} \\ \hline
\endhead
University Student & Primary & A registered undergraduate student who uses the forum, shares resources, and optionally joins study rooms; for UC-05, comments on or answers a forum post. \\ \hline
Post Author & Secondary & A student who receives notifications when comments or answers are added to their post. \\ \hline
Notification Service & Supporting System & Delivers in-app notifications for forum and room activity; for UC-05, delivers the forum notification to the post author. \\ \hline
\end{longtable}
\endgroup
\par\medskip

\hypertarget{sec:a3-relationships}{}
\addcontentsline{toc}{subsubsection}{A3. Relationships}
\subsubsection*{A3. Relationships}
\par\medskip

\hypertarget{sec:a3-1-associations}{}
\paragraph*{A3.1 Associations}
\par\medskip

\begingroup
\small
\renewcommand{\baselinestretch}{1.0}\selectfont
\renewcommand{\arraystretch}{1.15}
\setlength{\tabcolsep}{6pt}
\begin{longtable}{|>{\raggedright\arraybackslash}p{0.16\textwidth}|>{\raggedright\arraybackslash}p{0.22\textwidth}|>{\raggedright\arraybackslash}p{0.10\textwidth}|>{\raggedright\arraybackslash}p{0.44\textwidth}|}
\hline
\textbf{Actor} & \textbf{Use Case} & \textbf{Type} & \textbf{Description} \\ \hline
\endhead
University Student & UC-05 Comment and Answer & Primary & Comments on or answers a forum post. \\ \hline
Post Author & I-06 Send Forum Notification & Secondary & Receives notification when a comment or answer is added. \\ \hline
Notification Service & I-06 Send Forum Notification & Supporting & Delivers the forum notification to the post author. \\ \hline
\end{longtable}
\endgroup
\par\medskip

\hypertarget{sec:a3-2-generalization-and-specialization}{}
\paragraph*{A3.2 Generalization and Specialization}
\par\medskip

\begingroup
\small
\renewcommand{\baselinestretch}{1.0}\selectfont
\renewcommand{\arraystretch}{1.15}
\setlength{\tabcolsep}{6pt}
\begin{longtable}{|>{\raggedright\arraybackslash}p{0.10\textwidth}|>{\raggedright\arraybackslash}p{0.11\textwidth}|>{\raggedright\arraybackslash}p{0.12\textwidth}|>{\raggedright\arraybackslash}p{0.59\textwidth}|}
\hline
\textbf{Child} & \textbf{Parent} & \textbf{Type} & \textbf{Inherited Behavior or Association} \\ \hline
\endhead
Post Author & University Student & Actor generalization & Inherits the general student associations, including UC-05, and adds I-06 as notification recipient. \\ \hline
\end{longtable}
\endgroup
\par\medskip

\hypertarget{sec:a3-3-include}{}
\paragraph*{A3.3 Include}
\par\medskip

\begingroup
\small
\renewcommand{\baselinestretch}{1.0}\selectfont
\renewcommand{\arraystretch}{1.15}
\setlength{\tabcolsep}{6pt}
\begin{longtable}{|>{\raggedright\arraybackslash}p{0.21\textwidth}|>{\raggedright\arraybackslash}p{0.24\textwidth}|>{\raggedright\arraybackslash}p{0.47\textwidth}|}
\hline
\textbf{Base Use Case} & \textbf{Included Use Case} & \textbf{Meaning} \\ \hline
\endhead
UC-05 Comment and Answer & I-06 Send Forum Notification & Notify the post author about the new comment or answer. \\ \hline
\end{longtable}
\endgroup
\par\medskip

\hypertarget{sec:a3-4-extend}{}
\paragraph*{A3.4 Extend}
\par\medskip

\begingroup
\small
\renewcommand{\baselinestretch}{1.0}\selectfont
\renewcommand{\arraystretch}{1.15}
\setlength{\tabcolsep}{6pt}
\begin{longtable}{|>{\raggedright\arraybackslash}p{0.13\textwidth}|>{\raggedright\arraybackslash}p{0.17\textwidth}|>{\raggedright\arraybackslash}p{0.38\textwidth}|>{\raggedright\arraybackslash}p{0.22\textwidth}|}
\hline
\textbf{Extending Use Case} & \textbf{Base Use Case} & \textbf{Condition} & \textbf{Added Behavior} \\ \hline
\endhead
E-06 Nested Reply & UC-05 Comment and Answer & The student replies to an existing comment or answer. & Add a threaded nested response. \\ \hline
\end{longtable}
\endgroup
\par\medskip

\hypertarget{sec:a4-model-notes}{}
\addcontentsline{toc}{subsubsection}{A4. Model Notes}
\subsubsection*{A4. Model Notes}
\par\medskip

\begin{enumerate}
\item Every new comment or answer includes I-06 because the post author must be notified.
\item Nested reply is an optional extension of the base comment-or-answer behavior.
\item Post Author inherits the University Student association to UC-05 and adds the secondary notification-recipient association to I-06.
\end{enumerate}
\par\medskip

\hypertarget{sec:part-b-use-case-specification}{}
\addcontentsline{toc}{subsection}{Part B --- Use Case Specification}
\subsection*{Part B --- Use Case Specification}
\par\medskip

\hypertarget{sec:b1-brief-description}{}
\addcontentsline{toc}{subsubsection}{B1. Brief Description}
\subsubsection*{B1. Brief Description}
\par\medskip

A student contributes to an academic discussion by commenting on a post or answering a question, so that ideas are discussed and other students receive help. The post author is notified about each new comment or answer. UC-05 covers creation, nested replies, edit/delete of own responses, answer classification for question posts, discussion display, and notification behavior in the primary University Forum feature group.
\par\medskip

\hypertarget{sec:b2-preconditions}{}
\addcontentsline{toc}{subsubsection}{B2. Preconditions}
\subsubsection*{B2. Preconditions}
\par\medskip

\begin{enumerate}
\item The student is authenticated with an active session.
\item The student can view the selected post.
\item The post accepts comments or answers and has not been removed or locked.
\item The student is not restricted from commenting or answering.
\end{enumerate}
\par\medskip

\hypertarget{sec:b3-postconditions}{}
\addcontentsline{toc}{subsubsection}{B3. Postconditions}
\subsubsection*{B3. Postconditions}
\par\medskip

\begin{enumerate}
\item \textbf{Success:} The comment or answer is stored and associated with the post, its author, and its parent comment when applicable.
\item \textbf{Success:} The comment or answer is displayed in the post's discussion in chronological or threaded order.
\item \textbf{Success:} The post author receives a notification about the new comment or answer.
\item \textbf{Failure:} No comment or answer is stored, and the student receives a clear error message.
\end{enumerate}
\par\medskip

\hypertarget{sec:b4-main-success-scenario-basic-flow}{}
\addcontentsline{toc}{subsubsection}{B4. Main Success Scenario (Basic Flow)}
\subsubsection*{B4. Main Success Scenario (Basic Flow)}
\par\medskip

Unless otherwise stated, an alternative flow resumes at the next step of the basic flow.
\par\medskip

\begin{enumerate}
\item The student opens a post detail page.
\item The system displays the post's existing comments and answers, including the marked true response when present.
\item The student selects the comment or answer composer.
\item The student enters the response content.
\item The system identifies whether the response is a comment on a discussion post or an answer to a question post.
\item The student submits the response.
\item The frontend validates that the response is not empty and submits it to the backend.
\item The backend verifies authentication, post visibility, posting permission, content validity, and rate limits.
\item The backend stores the comment or answer, associates it with the post and the student, and updates the discussion structure.
\item The system updates the displayed thread and comment or answer count.
\item The notification service notifies the post author that a new comment or answer was added.
\item The student sees the published response in the thread.
\end{enumerate}
\par\medskip

\hypertarget{sec:b5-alternative-flows}{}
\addcontentsline{toc}{subsubsection}{B5. Alternative Flows}
\subsubsection*{B5. Alternative Flows}
\par\medskip

\hypertarget{sec:a1-nested-reply}{}
\paragraph*{A1. Nested reply}
\par\medskip

At basic flow step 4, the student replies to an existing comment or answer. The system stores the response with its parent and displays it in the threaded structure where applicable.
\par\medskip

\hypertarget{sec:a2-edit-own-response}{}
\paragraph*{A2. Edit own response}
\par\medskip

The student selects \textbf{Edit} on their own comment or answer, changes the content, and submits. The system validates and updates the existing response.
\par\medskip

\hypertarget{sec:a3-delete-own-response}{}
\paragraph*{A3. Delete own response}
\par\medskip

The student selects \textbf{Delete} on their own comment or answer and confirms. The system removes or hides the response according to the platform's deletion policy and updates the thread.
\par\medskip

\hypertarget{sec:a4-answer-a-question}{}
\paragraph*{A4. Answer a question}
\par\medskip

At basic flow step 5, the system classifies the response as an answer to the question post, making it eligible for true-response marking by the post author.
\par\medskip

\hypertarget{sec:a5-no-existing-discussion}{}
\paragraph*{A5. No existing discussion}
\par\medskip

At basic flow step 2, the system shows an empty discussion state and still presents the composer.
\par\medskip

\hypertarget{sec:b6-exception-flows}{}
\addcontentsline{toc}{subsubsection}{B6. Exception Flows}
\subsubsection*{B6. Exception Flows}
\par\medskip

\hypertarget{sec:e1-empty-or-invalid-response}{}
\paragraph*{E1. Empty or invalid response}
\par\medskip

The system keeps the composer open, shows validation feedback, and does not submit.
\par\medskip

\hypertarget{sec:e2-post-removed-locked-or-inaccessible}{}
\paragraph*{E2. Post removed, locked, or inaccessible}
\par\medskip

The system rejects the response and explains that the post is no longer available for discussion.
\par\medskip

\hypertarget{sec:e3-posting-restriction}{}
\paragraph*{E3. Posting restriction}
\par\medskip

A restricted or unauthorized student is denied commenting or answering.
\par\medskip

\hypertarget{sec:e4-rate-limit-reached}{}
\paragraph*{E4. Rate limit reached}
\par\medskip

The system rejects the response and asks the student to retry later.
\par\medskip

\hypertarget{sec:e5-notification-delivery-failure}{}
\paragraph*{E5. Notification delivery failure}
\par\medskip

The comment or answer remains published; the system records or retries the notification according to the notification policy.
\par\medskip

\hypertarget{sec:e6-server-or-database-failure}{}
\paragraph*{E6. Server or database failure}
\par\medskip

No response is stored, and the student receives a retry option with the entered content preserved where supported.
\par\medskip

\hypertarget{sec:b7-business-rules}{}
\addcontentsline{toc}{subsubsection}{B7. Business Rules}
\subsubsection*{B7. Business Rules}
\par\medskip

\begin{enumerate}
\item Comments are allowed on discussion posts; answers are allowed on question posts.
\item Nested replies are supported where applicable.
\item Students may edit or delete only their own comments and answers.
\item Comments and answers are displayed in chronological or threaded order.
\item The post author is notified when a new comment or answer is added.
\item Commenting and answering are subject to authentication, permission, and rate-limit checks.
\item Only one true response may be marked for a question post; that action is performed by the post author under US-20.
\end{enumerate}
\par\medskip

\hypertarget{sec:b8-special-requirements}{}
\addcontentsline{toc}{subsubsection}{B8. Special Requirements}
\subsubsection*{B8. Special Requirements}
\par\medskip

\begingroup
\small
\renewcommand{\baselinestretch}{1.0}\selectfont
\renewcommand{\arraystretch}{1.15}
\setlength{\tabcolsep}{6pt}
\begin{longtable}{|>{\raggedright\arraybackslash}p{0.10\textwidth}|>{\raggedright\arraybackslash}p{0.83\textwidth}|}
\hline
\textbf{ID} & \textbf{Requirement} \\ \hline
\endhead
NFR-PERF-4 & A new comment or answer must be reflected in the relevant forum view within 2 seconds. \\ \hline
NFR-PERF-6 & Forum notifications must be delivered with a latency of 2 seconds or less under normal operating conditions. \\ \hline
NFR-SEC-9 & Commenting and answering must be rate-limited to reduce spam and abuse. \\ \hline
Design constraint & A notification delivery failure must not roll back an already published comment or answer; the notification may be recorded or retried. \\ \hline
\end{longtable}
\endgroup
\par\medskip

\hypertarget{sec:b9-extension-points}{}
\addcontentsline{toc}{subsubsection}{B9. Extension Points}
\subsubsection*{B9. Extension Points}
\par\medskip

\begingroup
\small
\renewcommand{\baselinestretch}{1.0}\selectfont
\renewcommand{\arraystretch}{1.15}
\setlength{\tabcolsep}{6pt}
\begin{longtable}{|>{\raggedright\arraybackslash}p{0.32\textwidth}|>{\raggedright\arraybackslash}p{0.61\textwidth}|}
\hline
\textbf{Extension Point} & \textbf{Location in Flow of Events} \\ \hline
\endhead
Nested reply & Basic flow step 4; alternative A1 \\ \hline
Edit or delete own response & After the response is displayed; alternatives A2-A3 \\ \hline
Answer classification & Basic flow step 5; alternative A4 \\ \hline
Empty discussion & Basic flow step 2; alternative A5 \\ \hline
Notification delivery & Basic flow step 11; exception E5 \\ \hline
Rate-limit recovery & Basic flow step 8; exception E4 \\ \hline
\end{longtable}
\endgroup
\par\medskip

\hypertarget{sec:b10-interface-and-data-notes}{}
\addcontentsline{toc}{subsubsection}{B10. Interface and Data Notes}
\subsubsection*{B10. Interface and Data Notes}
\par\medskip

\begin{itemize}
\item \textbf{Screen flow.} \texttt{screen\_flow.md}: Home $\rightarrow$ Forum Group $\rightarrow$ University $\rightarrow$ Major $\rightarrow$ Course $\rightarrow$ Post, with Trending $\rightarrow$ Post and Create Post $\rightarrow$ Post. \texttt{feature\_flow.md}: List Post $\rightarrow$ Post $\rightarrow$ Comment / Solved / Vote / Save / Report / Moderate; Post $\rightarrow$ Comment $\rightarrow$ True Response; Comment $\rightarrow$ Vote. \texttt{post\_management.md} recommends the UI route \texttt{/forum/posts/:postId} and the detail components \texttt{PostDetailPage}, \texttt{CommentList}, \texttt{CommentComposer}, \texttt{TrueResponsePanel}, \texttt{PostSolvedBadge}, and \texttt{PermissionBasedActionBar}; \texttt{API.md} uses \texttt{/forum/posts/:post\_id}.
\item \textbf{Detail load.} \texttt{GET /forum/posts/:post\_id} returns \texttt{post}, \texttt{comments}, \texttt{true\_response}, \texttt{user\_vote}, and \texttt{saved} with \texttt{200}; failure states are \texttt{404 POST\_NOT\_FOUND} (\texttt{48001}), \texttt{401}, and \texttt{500 (DATABASE\_ERROR | UNKNOWN\_ERROR)}. Backend function \texttt{studitopia.get\_forum\_post(ppost\_id: uuid)} also lists \texttt{FORBIDDEN} (\texttt{48010}).
\item \textbf{Composer/create.} \texttt{CommentComposer} is a textarea plus submit control for comments, answers, and replies (\texttt{FR-18.1}--\texttt{FR-18.3}, \texttt{AC-18.1}--\texttt{AC-18.3}). Creation uses \texttt{POST /forum/posts/:post\_id/comments} with \texttt{\{content, parent\_comment\_id?\}}; success is \texttt{201 \{comment\_id, created\_at\}}. Errors are \texttt{400 INVALID\_POST\_CONTENT} (\texttt{48004}), \texttt{404 POST\_NOT\_FOUND} (\texttt{48001}) or \texttt{COMMENT\_NOT\_FOUND} (\texttt{48002}), \texttt{401}, and \texttt{500}. Database function \texttt{studitopia.create\_forum\_comment(pauthor\_id, ppost\_id, pcontent, pparent\_comment\_id)} additionally lists \texttt{UNKNOWN\_ERROR} (\texttt{45000}) and \texttt{FORBIDDEN} (\texttt{48010}).
\item \textbf{Edit/delete rules.} \texttt{CommentActionMenu}, \texttt{EditCommentDialog}, and \texttt{DeleteCommentDialog} expose \textbf{Edit} and \textbf{Delete} only for the current user's own comment or answer (\texttt{FR-18.4}, \texttt{AC-18.4}, business rule 3). \texttt{post\_management.md} requires backend enforcement but states that no comment edit/delete API is specified in \texttt{API.md}; no endpoint is invented here. \texttt{UC-05} A3 removes or hides the response according to platform deletion policy.
\item \textbf{Notification rules.} Every new response includes I-06 Send Forum Notification (\texttt{FR-18.6}, \texttt{FR-38.1}, UC-05 step 11). \texttt{settings.md} surfaces \texttt{ForumNotificationToggles} for comments, answers, votes, true-response marks, and mentions; \texttt{AudioAlertToggle}; \texttt{NotificationCenterLink}; \texttt{MarkAllReadButton}; and \texttt{SaveNotificationSettingsButton}, with \texttt{NotificationPreference} fields \texttt{forumEvents}, \texttt{socialEvents}, \texttt{roomInvitationsEnabled}, \texttt{audioAlertsEnabled}. No notification endpoint is defined in \texttt{API.md}; \texttt{E5} records or retries delivery and must not roll back the published response.
\item \textbf{True response.} \texttt{PATCH /forum/posts/:post\_id/true-response} with \texttt{\{comment\_id: string\}} returns \texttt{200 \{solved: boolean\}}; errors are \texttt{404 POST\_NOT\_FOUND} (\texttt{48001}) or \texttt{COMMENT\_NOT\_FOUND} (\texttt{48002}), \texttt{403 NOT\_POST\_OWNER} (\texttt{48007}) or \texttt{FORBIDDEN} (\texttt{48010}), \texttt{401}, and \texttt{500}. Backend function \texttt{studitopia.mark\_true\_response(puser\_id, ppost\_id, pcomment\_id)} lists the same codes. \texttt{PostSolvedBadge} shows \texttt{Solved} / \texttt{Unsolved}; \texttt{TrueResponsePanel} shows the accepted answer or comment prominently (\texttt{FR-20.1}--\texttt{FR-20.5}).
\item \textbf{Save.} \texttt{POST /forum/posts/:post\_id/save} returns \texttt{200 \{saved: boolean\}}; errors are \texttt{404 POST\_NOT\_FOUND} (\texttt{48001}), \texttt{401}, and \texttt{500}. Database function \texttt{studitopia.save\_forum\_post(puser\_id, ppost\_id)} also lists \texttt{FORBIDDEN} (\texttt{48010}). \texttt{SavePostButton} shows \textbf{Save}/\textbf{Saved} state (\texttt{FR-21.1}--\texttt{FR-21.4}).
\item \textbf{Report and moderation.} \texttt{POST /forum/reports} with \texttt{\{target\_type: 'post' | 'comment', target\_id, category, description?\}} returns \texttt{201 \{report\_id\}}; errors are \texttt{400 INVALID\_REPORT\_CATEGORY} (\texttt{48009}), \texttt{404 POST\_NOT\_FOUND} (\texttt{48001}) or \texttt{COMMENT\_NOT\_FOUND} (\texttt{48002}), \texttt{401}, and \texttt{500}. Database function \texttt{studitopia.report\_forum\_content(preporter\_id, ptarget\_type, ptarget\_id, pcategory, pdescription)} also lists \texttt{FORBIDDEN} (\texttt{48010}). Reports route to the moderation queue; duplicate-report prevention is stated in \texttt{API.md} but has no specified error code. \texttt{moderation\_queue.md} has no queue/action/history endpoints and notes that moderation actions are restricted to authorized moderators/administrators (\texttt{NFR-SEC-10}).
\item \textbf{Vote.} \texttt{POST /forum/votes} with \texttt{\{target\_type: 'post' | 'comment', target\_id, vote: 1 | 0 | -1\}} returns \texttt{200 \{vote\_score\}}; errors are \texttt{400 INVALID\_VOTE} (\texttt{48005}), \texttt{409 DUPLICATED\_VOTE} (\texttt{48006}), \texttt{404 POST\_NOT\_FOUND} (\texttt{48001}) or \texttt{COMMENT\_NOT\_FOUND} (\texttt{48002}), \texttt{401}, and \texttt{500}. Database function \texttt{studitopia.vote\_forum\_content(puser\_id, ptarget\_type, ptarget\_id, pvote)} also lists \texttt{FORBIDDEN} (\texttt{48010}). One active vote per target may be changed or removed (\texttt{FR-19.1}--\texttt{FR-19.5}, \texttt{AC-19.3}, \texttt{AC-19.4}).
\item \textbf{Other forum codes.} \texttt{response\_code.md} defines \texttt{48000 UNKNOWN FORUM ERROR}, \texttt{48003 INVALID\_POST\_TYPE}, \texttt{48008 ALREADY\_TRUE\_RESPONSE}, and \texttt{48010 FORBIDDEN}. \texttt{48008} is defined but is not listed in the \texttt{PATCH /forum/posts/:post\_id/true-response} response contract.
\item \textbf{Privacy and permission rules.} \texttt{NFR-SEC-1} requires authentication for protected resources; \texttt{NFR-SEC-4} restricts private forum content to authorized users; \texttt{NFR-SEC-6} requires server-side authorization. \texttt{NFR-PRIV-7} applies privacy settings consistently to forum posts, comments, member lists, and search results. \texttt{FR-17.4}/\texttt{AC-17.4} display only permitted actions, and \texttt{FR-39.5} requires backend prevention of disallowed interactions. \texttt{NFR-SEC-9} rate-limits commenting; \texttt{NFR-USE-6} places report/moderation actions within at most 2 interactions from a post or comment. \texttt{friend.md} and \texttt{friend\_profile.md} show friend posts/comments only under privacy rules and do not define a comment API.
\end{itemize}
\par\medskip

\hypertarget{sec:b11-source-alignment-notes}{}
\addcontentsline{toc}{subsubsection}{B11. Source Alignment Notes}
\subsubsection*{B11. Source Alignment Notes}
\par\medskip

\begin{itemize}
\item The focused model file and \texttt{use\_case\_diagrams.md} contain identical UC-05 diagram code; the model index calls the Post Author link "actor inheritance" while the focused model labels it "generalization". This export uses the focused model's terminology.
\item \texttt{use\_cases\_description.md} and \texttt{PA3-Group05/use\_case.md} agree on UC-05's goal, 12-step basic flow, A1--A5, E1--E6, business rules, related requirements, special requirements, and extension points.
\item \texttt{API.md} exposes one comment-creation endpoint for comments, answers, and nested replies (\texttt{parent\_comment\_id?}) but no answer-specific endpoint, no comment edit/delete endpoint, and no notification endpoint; edit/delete remains an \texttt{FR-18.4} UI requirement with unspecified backend contract.
\item \texttt{UC-05} states that answers are classified separately from comments, and \texttt{FR-18.2} requires answer creation, but \texttt{API.md} does not define a separate answer type or field for that classification.
\item \texttt{FR-18.6}/UC-05 state that the post author is notified when a new response is added, while \texttt{FR-38.5} allows users to configure which forum events trigger alerts; no source specifies how mandatory notification behavior interacts with user opt-outs.
\item \texttt{ALREADY\_TRUE\_RESPONSE} (\texttt{48008}) exists in \texttt{response\_code.md}, but \texttt{PATCH /forum/posts/:post\_id/true-response} lists only \texttt{NOT\_POST\_OWNER} (\texttt{48007}) and \texttt{FORBIDDEN} (\texttt{48010}) under \texttt{403}; the one-true-response rule comes from business rule 7 and \texttt{FR-20.1}/\texttt{FR-20.4}.
\item \texttt{screen\_flow.md} and \texttt{feature\_flow.md} define navigation but not exact routes; \texttt{post\_management.md} marks \texttt{/forum/posts/:postId} as recommended, while \texttt{API.md} uses \texttt{/forum/posts/:post\_id}.
\item \texttt{E5} and the PA3 design constraint specify that notification failure must not roll back a published response, but the \texttt{500} response contract does not describe partial notification state.
\item \texttt{friend.md} and \texttt{friend\_profile.md} explicitly state that there is no friend/profile API contract; they contribute only privacy-scoped visibility of friend posts/comments (\texttt{NFR-PRIV-7}), not additional forum endpoints.
\end{itemize}
\par\medskip

\par\medskip
\hypertarget{sec:ui-prototype-specimens}{}
\addcontentsline{toc}{subsection}{UI Prototype Specimens}
\subsection*{UI Prototype Specimens}
\par\medskip
\begin{figure}[H]
\centering
\includegraphics[width=0.75\textwidth,height=0.42\textheight,keepaspectratio]{"../../prototype/dijkstra_question_desktop.png"}
\par\vspace{2pt}
{\footnotesize \textbf{UI Prototype:} Dijkstra question - Desktop}
\end{figure}
\par\smallskip

\pagebreak
\hypertarget{sec:uc-17-vote-on-posts-and-comments}{}
\addcontentsline{toc}{section}{UC-17 --- Vote on Posts and Comments}
\section*{UC-17 --- Vote on Posts and Comments}
\par\medskip

\textbf{Version:} \texttt{1.0}
\par\medskip

\textbf{Performed by:} [Nguyen Le Quoc Huy] | \textbf{Edited by:} [Nguyen Le Quoc Huy]
\par\medskip

\begingroup
\small
\renewcommand{\baselinestretch}{1.0}\selectfont
\renewcommand{\arraystretch}{1.15}
\setlength{\tabcolsep}{6pt}
\begin{longtable}{|>{\raggedright\arraybackslash}p{0.09\textwidth}|>{\raggedright\arraybackslash}p{0.84\textwidth}|}
\hline
\textbf{Field} & \textbf{Value} \\ \hline
\endhead
Use Case ID & UC-17 \\ \hline
Name & Vote on Posts and Comments \\ \hline
Primary Actor & University Student \\ \hline
Supporting Actors & None \\ \hline
Feature Group & Forum Interaction \\ \hline
Priority & Not explicitly stated \\ \hline
Related User Stories & US-19 \\ \hline
Related Requirements & FR-13.2, FR-13.4, FR-17.4, FR-19.1 to FR-19.5, FR-38.1; NFR-PERF-2, NFR-PERF-3, NFR-PERF-4, NFR-PERF-6, NFR-SEC-1, NFR-SEC-6, NFR-SEC-9, NFR-PRIV-7 \\ \hline
Source Documents & \texttt{proposals/userstory.md} (section 6.2); \texttt{proposals/functional\_requirement.md} (US-19); \texttt{proposals/list\_post.md}; \texttt{proposals/post\_management.md}; \texttt{proposals/create\_post.md}; \texttt{proposals/API.md}; \texttt{proposals/response\_code.md}; \texttt{proposals/non\_functional\_requirement.md}; \texttt{proposals/feature\_flow.md}; \texttt{proposals/screen\_flow.md} \\ \hline
\end{longtable}
\endgroup
\par\medskip

\hypertarget{sec:part-a-use-case-model}{}
\addcontentsline{toc}{subsection}{Part A --- Use Case Model}
\subsection*{Part A --- Use Case Model}
\par\medskip

\hypertarget{sec:a1-focused-use-case-diagram}{}
\addcontentsline{toc}{subsubsection}{A1. Focused Use Case Diagram}
\subsubsection*{A1. Focused Use Case Diagram}
\par\medskip

\begin{center}
\includegraphics[width=0.85\textwidth,height=0.75\textheight,keepaspectratio]{"figures/UC-17-vote-posts-comments_1.png"}
\end{center}
\par\medskip

\textbf{Diagram notation.} Person shape = human actor; cloud = supporting system; blue = selected use case; solid line = association; labelled solid arrow = generalization; dashed arrow = include or extend.
\par\medskip

\hypertarget{sec:a2-actor-register}{}
\addcontentsline{toc}{subsubsection}{A2. Actor Register}
\subsubsection*{A2. Actor Register}
\par\medskip

\begingroup
\small
\renewcommand{\baselinestretch}{1.0}\selectfont
\renewcommand{\arraystretch}{1.15}
\setlength{\tabcolsep}{6pt}
\begin{longtable}{|>{\raggedright\arraybackslash}p{0.11\textwidth}|>{\raggedright\arraybackslash}p{0.11\textwidth}|>{\raggedright\arraybackslash}p{0.70\textwidth}|}
\hline
\textbf{Actor} & \textbf{Type} & \textbf{Description} \\ \hline
\endhead
University Student & Primary human actor & A student who upvotes or downvotes a post, comment, or answer so that useful academic content is easier to discover. \\ \hline
\end{longtable}
\endgroup
\par\medskip

\hypertarget{sec:a3-relationships}{}
\addcontentsline{toc}{subsubsection}{A3. Relationships}
\subsubsection*{A3. Relationships}
\par\medskip

\hypertarget{sec:a3-1-associations}{}
\paragraph*{A3.1 Associations}
\par\medskip

\begingroup
\small
\renewcommand{\baselinestretch}{1.0}\selectfont
\renewcommand{\arraystretch}{1.15}
\setlength{\tabcolsep}{6pt}
\begin{longtable}{|>{\raggedright\arraybackslash}p{0.12\textwidth}|>{\raggedright\arraybackslash}p{0.23\textwidth}|>{\raggedright\arraybackslash}p{0.10\textwidth}|>{\raggedright\arraybackslash}p{0.46\textwidth}|}
\hline
\textbf{Actor} & \textbf{Use Case} & \textbf{Type} & \textbf{Description} \\ \hline
\endhead
University Student & UC-17 Vote on Posts and Comments & Primary & Casts, changes, or removes a vote on an eligible post or comment. \\ \hline
\end{longtable}
\endgroup
\par\medskip

\hypertarget{sec:a3-2-generalization-and-specialization}{}
\paragraph*{A3.2 Generalization and Specialization}
\par\medskip

\begingroup
\small
\renewcommand{\baselinestretch}{1.0}\selectfont
\renewcommand{\arraystretch}{1.15}
\setlength{\tabcolsep}{6pt}
\begin{longtable}{|>{\raggedright\arraybackslash}p{0.09\textwidth}|>{\raggedright\arraybackslash}p{0.09\textwidth}|>{\raggedright\arraybackslash}p{0.09\textwidth}|>{\raggedright\arraybackslash}p{0.65\textwidth}|}
\hline
\textbf{Child} & \textbf{Parent} & \textbf{Type} & \textbf{Inherited Behavior or Association} \\ \hline
\endhead
None & None & None & No actor generalization or specialization is named for US-19 in the specified proposals. \\ \hline
\end{longtable}
\endgroup
\par\medskip

\hypertarget{sec:a3-3-include}{}
\paragraph*{A3.3 Include}
\par\medskip

\begingroup
\small
\renewcommand{\baselinestretch}{1.0}\selectfont
\renewcommand{\arraystretch}{1.15}
\setlength{\tabcolsep}{6pt}
\begin{longtable}{|>{\raggedright\arraybackslash}p{0.10\textwidth}|>{\raggedright\arraybackslash}p{0.11\textwidth}|>{\raggedright\arraybackslash}p{0.71\textwidth}|}
\hline
\textbf{Base Use Case} & \textbf{Included Use Case} & \textbf{Meaning} \\ \hline
\endhead
None & None & No distinct included use case is named for US-19 in the specified proposals; no new I-xx ID is introduced. \\ \hline
\end{longtable}
\endgroup
\par\medskip

\hypertarget{sec:a3-4-extend}{}
\paragraph*{A3.4 Extend}
\par\medskip

\begingroup
\small
\renewcommand{\baselinestretch}{1.0}\selectfont
\renewcommand{\arraystretch}{1.15}
\setlength{\tabcolsep}{6pt}
\begin{longtable}{|>{\raggedright\arraybackslash}p{0.10\textwidth}|>{\raggedright\arraybackslash}p{0.09\textwidth}|>{\raggedright\arraybackslash}p{0.09\textwidth}|>{\raggedright\arraybackslash}p{0.62\textwidth}|}
\hline
\textbf{Extending Use Case} & \textbf{Base Use Case} & \textbf{Condition} & \textbf{Added Behavior} \\ \hline
\endhead
None & None & None & No distinct extending use case is named for US-19 in the specified proposals; no new E-xx ID is introduced. \\ \hline
\end{longtable}
\endgroup
\par\medskip

\hypertarget{sec:a4-model-notes}{}
\addcontentsline{toc}{subsubsection}{A4. Model Notes}
\subsubsection*{A4. Model Notes}
\par\medskip

\begin{enumerate}
\item UC-17 covers both target types named in \texttt{API.md}: \texttt{post} and \texttt{comment}.
\item Comments and answers share the \texttt{comment} target type in the API; no separate answer vote endpoint is named.
\item Upvote, downvote, change, and remove are values of one vote action, not separate use cases.
\item Notification of vote activity is a related cross-cutting behavior in \texttt{FR-38.1}, not a source-named included use case; priority and a dedicated vote UI route are not specified in proposals.
\end{enumerate}
\par\medskip

\hypertarget{sec:part-b-use-case-specification}{}
\addcontentsline{toc}{subsection}{Part B --- Use Case Specification}
\subsection*{Part B --- Use Case Specification}
\par\medskip

\hypertarget{sec:b1-brief-description}{}
\addcontentsline{toc}{subsubsection}{B1. Brief Description}
\subsubsection*{B1. Brief Description}
\par\medskip

A university student votes on a post or on a comment or answer so that useful academic content is easier to discover. The student can upvote, downvote, change a vote, or remove a vote, with at most one active vote per target. The recorded vote updates the vote score and ranking, and vote controls appear where the user has permission to view and interact with the target.
\par\medskip

\hypertarget{sec:b2-preconditions}{}
\addcontentsline{toc}{subsubsection}{B2. Preconditions}
\subsubsection*{B2. Preconditions}
\par\medskip

\begin{enumerate}
\item The student is authenticated with an active session.
\item The student can view the selected post or comment.
\item The target is a \texttt{post} or \texttt{comment} and still exists.
\item The student is permitted to vote under privacy, moderation, and authorization rules.
\end{enumerate}
\par\medskip

\hypertarget{sec:b3-postconditions}{}
\addcontentsline{toc}{subsubsection}{B3. Postconditions}
\subsubsection*{B3. Postconditions}
\par\medskip

\begin{enumerate}
\item \textbf{Success:} The vote is recorded as \texttt{1}, \texttt{0}, or \texttt{-1} for the target and requesting user.
\item \textbf{Success:} The vote score and ranking are updated after the vote is recorded.
\item \textbf{Success:} The user's vote state is reflected in the forum view.
\item \textbf{Failure:} The vote is not recorded, the previous vote state remains authoritative, and the student receives a clear error.
\end{enumerate}
\par\medskip

\hypertarget{sec:b4-main-success-scenario-basic-flow}{}
\addcontentsline{toc}{subsubsection}{B4. Main Success Scenario (Basic Flow)}
\subsubsection*{B4. Main Success Scenario (Basic Flow)}
\par\medskip

\begin{enumerate}
\item The student opens a forum post list or post detail page.
\item The system displays the target's current vote score and the student's current vote state.
\item The student selects upvote or downvote on a post, comment, or answer.
\item The frontend sends \texttt{target\_type}, \texttt{target\_id}, and the selected \texttt{vote} value.
\item The backend validates the target type, target existence, vote value, permissions, and rate limits.
\item The system records or replaces the student's vote for that target.
\item The system returns the updated \texttt{vote\_score}.
\item The system refreshes the vote score, ranking, and vote controls in the forum view.
\end{enumerate}
\par\medskip

\hypertarget{sec:b5-alternative-flows}{}
\addcontentsline{toc}{subsubsection}{B5. Alternative Flows}
\subsubsection*{B5. Alternative Flows}
\par\medskip

Unless otherwise stated, an alternative flow resumes at the next step of the basic flow.
\par\medskip

\hypertarget{sec:a1-vote-on-a-comment-or-answer}{}
\paragraph*{A1. Vote on a comment or answer}
\par\medskip

At basic flow step 3, the student votes on a comment or answer. The request uses \texttt{target\_type: 'comment'} with the comment or answer identifier.
\par\medskip

\hypertarget{sec:a2-change-an-existing-vote}{}
\paragraph*{A2. Change an existing vote}
\par\medskip

At basic flow step 3, the student selects the opposite vote direction. The existing vote is replaced for the same target, preserving one active vote per user and target.
\par\medskip

\hypertarget{sec:a3-remove-a-vote}{}
\paragraph*{A3. Remove a vote}
\par\medskip

At basic flow step 3, the student removes the current vote by submitting \texttt{vote: 0}. The score is recalculated without an active vote for that target.
\par\medskip

\hypertarget{sec:a4-vote-from-list-or-detail}{}
\paragraph*{A4. Vote from list or detail}
\par\medskip

At basic flow step 1, the student votes from the post list or from post detail where the corresponding control is displayed.
\par\medskip

\hypertarget{sec:b6-exception-flows}{}
\addcontentsline{toc}{subsubsection}{B6. Exception Flows}
\subsubsection*{B6. Exception Flows}
\par\medskip

\hypertarget{sec:e1-invalid-vote-value-or-target-type}{}
\paragraph*{E1. Invalid vote value or target type}
\par\medskip

The system returns \texttt{400} with \texttt{INVALID\_VOTE} (\texttt{48005}) and does not change the score.
\par\medskip

\hypertarget{sec:e2-duplicate-vote-conflict}{}
\paragraph*{E2. Duplicate vote conflict}
\par\medskip

The system returns \texttt{409} with \texttt{DUPLICATED\_VOTE} (\texttt{48006}). The UI reflects the current vote and offers change or removal rather than creating a second active vote.
\par\medskip

\hypertarget{sec:e3-target-not-found}{}
\paragraph*{E3. Target not found}
\par\medskip

The system returns \texttt{404} with \texttt{POST\_NOT\_FOUND} (\texttt{48001}) or \texttt{COMMENT\_NOT\_FOUND} (\texttt{48002}) and refreshes the unavailable view.
\par\medskip

\hypertarget{sec:e4-unauthorized-or-forbidden-vote}{}
\paragraph*{E4. Unauthorized or forbidden vote}
\par\medskip

The system returns \texttt{401}, or \texttt{403} with \texttt{FORBIDDEN} (\texttt{48010}), and hides or disables the action for unauthorized users.
\par\medskip

\hypertarget{sec:e5-rate-limit-reached}{}
\paragraph*{E5. Rate limit reached}
\par\medskip

The system rejects the vote and asks the student to retry later. \texttt{NFR-SEC-9} requires vote rate limiting, but no rate-limit status or error code is specified in proposals.
\par\medskip

\hypertarget{sec:e6-server-or-database-failure}{}
\paragraph*{E6. Server or database failure}
\par\medskip

The system returns \texttt{500} with \texttt{DATABASE\_ERROR} or \texttt{UNKNOWN\_ERROR} (\texttt{45000}), leaves the previous vote state authoritative, and offers retry where supported.
\par\medskip

\hypertarget{sec:b7-business-rules}{}
\addcontentsline{toc}{subsubsection}{B7. Business Rules}
\subsubsection*{B7. Business Rules}
\par\medskip

\begin{enumerate}
\item \texttt{AC-19.1} / \texttt{FR-19.1}: Users can upvote or downvote a post.
\item \texttt{AC-19.2} / \texttt{FR-19.2}: Users can upvote or downvote a comment or answer.
\item \texttt{AC-19.3} / \texttt{FR-19.3}: A user cannot vote more than once on the same post or comment.
\item \texttt{AC-19.4} / \texttt{FR-19.4}: Users can change or remove their vote.
\item \texttt{AC-19.5} / \texttt{FR-19.5}: The vote score and ranking update after a vote is recorded.
\item The API vote value is \texttt{1 | 0 | -1}, where \texttt{0} is removal; a vote must target \texttt{post} or \texttt{comment}.
\item Authentication, server-side authorization, privacy, and moderation rules determine whether vote controls are available (\texttt{NFR-SEC-1}, \texttt{NFR-SEC-6}, \texttt{NFR-PRIV-7}).
\item Voting is rate-limited to reduce spam and abuse (\texttt{NFR-SEC-9}).
\end{enumerate}
\par\medskip

\hypertarget{sec:b8-special-requirements}{}
\addcontentsline{toc}{subsubsection}{B8. Special Requirements}
\subsubsection*{B8. Special Requirements}
\par\medskip

\begingroup
\small
\renewcommand{\baselinestretch}{1.0}\selectfont
\renewcommand{\arraystretch}{1.15}
\setlength{\tabcolsep}{6pt}
\begin{longtable}{|>{\raggedright\arraybackslash}p{0.16\textwidth}|>{\raggedright\arraybackslash}p{0.77\textwidth}|}
\hline
\textbf{ID} & \textbf{Requirement} \\ \hline
\endhead
NFR-PERF-2, NFR-PERF-3 & The post list must load within 2 seconds, and post detail and comments must display within 2 seconds. \\ \hline
NFR-PERF-4 & A new vote must be reflected in the relevant forum view within 2 seconds. \\ \hline
NFR-PERF-6 & Forum notifications, including vote activity, must be delivered within 2 seconds under normal conditions. \\ \hline
NFR-SEC-1 & Authentication is required before accessing protected interactions. \\ \hline
NFR-SEC-6 & Server-side authorization checks must prevent unauthorized votes. \\ \hline
NFR-SEC-9 & Voting must be rate-limited to reduce spam and abuse. \\ \hline
NFR-PRIV-7 & Privacy settings apply consistently to forum posts, comments, member lists, and search results. \\ \hline
\end{longtable}
\endgroup
\par\medskip

\hypertarget{sec:b9-extension-points}{}
\addcontentsline{toc}{subsubsection}{B9. Extension Points}
\subsubsection*{B9. Extension Points}
\par\medskip

\begingroup
\small
\renewcommand{\baselinestretch}{1.0}\selectfont
\renewcommand{\arraystretch}{1.15}
\setlength{\tabcolsep}{6pt}
\begin{longtable}{|>{\raggedright\arraybackslash}p{0.32\textwidth}|>{\raggedright\arraybackslash}p{0.61\textwidth}|}
\hline
\textbf{Extension Point} & \textbf{Location in Flow of Events} \\ \hline
\endhead
Select vote target & Basic flow steps 1-3; alternatives A1 and A4 \\ \hline
Change or remove vote & Basic flow step 3; alternatives A2-A3 \\ \hline
Permission filtering & Basic flow step 2; exception E4 \\ \hline
Duplicate-vote recovery & Basic flow step 6; exception E2 \\ \hline
Rate-limit recovery & Basic flow step 5; exception E5 \\ \hline
\end{longtable}
\endgroup
\par\medskip

\hypertarget{sec:b10-interface-and-data-notes}{}
\addcontentsline{toc}{subsubsection}{B10. Interface and Data Notes}
\subsubsection*{B10. Interface and Data Notes}
\par\medskip

\begin{itemize}
\item \textbf{Screen/flow path:} \texttt{screen\_flow.md} gives \texttt{Forum -$>$ University -$>$ Major -$>$ Course -$>$ Post}, with \texttt{Trending -$>$ Post}. \texttt{feature\_flow.md} gives \texttt{List Post -$>$ Post}, \texttt{Post -$>$ Upvote / Downvote}, and \texttt{Comment -$>$ Vote}.
\item \textbf{UI components:} \texttt{list\_post.md} recommends \texttt{VoteUpButton}, \texttt{VoteDownButton}, \texttt{PostVoteScoreText}, and \texttt{EmptyPostListState}; \texttt{post\_management.md} recommends \texttt{VoteControl} with current score and user vote state.
\item \textbf{Post list and detail fields:} \texttt{GET /forum/posts} returns \texttt{post\_id}, \texttt{title}, \texttt{post\_type}, \texttt{author}, \texttt{vote\_score}, \texttt{comment\_count}, \texttt{solved}, and \texttt{created\_at} with \texttt{200}; failures are \texttt{400}, \texttt{401}, and \texttt{500}. \texttt{GET /forum/posts/:post\_id} returns \texttt{post}, \texttt{comments}, \texttt{true\_response}, \texttt{user\_vote: -1 | 0 | 1}, and \texttt{saved} with \texttt{200}; failures are \texttt{404 POST\_NOT\_FOUND} (\texttt{48001}), \texttt{401}, and \texttt{500 (DATABASE\_ERROR | UNKNOWN\_ERROR)}, with \texttt{FORBIDDEN: 48010} also declared by \texttt{studitopia.get\_forum\_post}.
\item \textbf{API endpoint and failures:} \texttt{POST /forum/votes} with body \texttt{\{ target\_type: 'post' | 'comment', target\_id: string, vote: 1 | 0 | -1 \}} returns \texttt{200} with \texttt{data: \{ vote\_score: number \}}. Failures are \texttt{400} with \texttt{(INVALID\_VOTE)[]} (\texttt{48005}), \texttt{409} with \texttt{(DUPLICATED\_VOTE)[]} (\texttt{48006}), \texttt{404} with \texttt{(POST\_NOT\_FOUND | COMMENT\_NOT\_FOUND)[]} (\texttt{48001}, \texttt{48002}), \texttt{401}, and \texttt{500} with \texttt{(DATABASE\_ERROR | UNKNOWN\_ERROR)[]}. \texttt{studitopia.vote\_forum\_content(puser\_id: uuid, ptarget\_type: string, ptarget\_id: uuid, pvote: integer)} returns \texttt{\{vote\_score: integer\}} and declares \texttt{UNKNOWN\_ERROR: 45000}, \texttt{POST\_NOT\_FOUND: 48001}, \texttt{COMMENT\_NOT\_FOUND: 48002}, \texttt{INVALID\_VOTE: 48005}, \texttt{DUPLICATED\_VOTE: 48006}, and \texttt{FORBIDDEN: 48010}.
\item \textbf{Vote state and notifications:} \texttt{list\_post.md} recommends \texttt{user\_vote} as the current user's \texttt{-1 | 0 | 1} state; \texttt{post\_management.md} says change/removal uses the same endpoint with \texttt{vote: 0}. \texttt{FR-38.1} and \texttt{NFR-PERF-6} include votes in forum activity notifications, but no vote-specific notification endpoint is defined in \texttt{API.md}.
\end{itemize}
\par\medskip

\hypertarget{sec:b11-source-alignment-notes}{}
\addcontentsline{toc}{subsubsection}{B11. Source Alignment Notes}
\subsubsection*{B11. Source Alignment Notes}
\par\medskip

\begin{itemize}
\item UC-04 lists voting as a downstream action after a Question or General Discussion post is published; UC-17 covers the vote action itself and does not repeat post-creation steps. UC-05 covers comment and answer creation and records vote context as related; UC-17 votes on those comments and answers without adding comments.
\item \texttt{AC-19.3} prevents more than one vote per target, \texttt{AC-19.4} allows changing or removing a vote, and \texttt{API.md} exposes both a \texttt{409 DUPLICATED\_VOTE} contract and \texttt{vote: 0} removal; \texttt{post\_management.md} resolves this at the UI layer by reflecting current state and offering change/removal.
\item \texttt{functional\_requirement.md} and \texttt{list\_post.md} refer to ranking updates, but \texttt{API.md} returns only \texttt{vote\_score} from the vote endpoint; no ranking response field is specified. \texttt{FR-38.1} includes votes in forum notifications, but no notification payload or endpoint is specified in the read sources.
\item Route names and component names in \texttt{list\_post.md} and \texttt{post\_management.md} are recommendations; \texttt{API.md} defines the exact endpoint and codes used above. No distinct included, extending, or specialized use case is named for US-19, so no new I-xx/E-xx/G-xx ID is introduced.
\end{itemize}
\par\medskip

\par\medskip
\hypertarget{sec:ui-prototype-specimens}{}
\addcontentsline{toc}{subsection}{UI Prototype Specimens}
\subsection*{UI Prototype Specimens}
\par\medskip
\begin{figure}[H]
\centering
\includegraphics[width=0.75\textwidth,height=0.42\textheight,keepaspectratio]{"../../prototype/forum_feed_interaction_states_desktop.png"}
\par\vspace{2pt}
{\footnotesize \textbf{UI Prototype:} Forum Feed - Interaction states - Desktop}
\end{figure}
\par\smallskip

\pagebreak
\hypertarget{sec:uc-18-mark-true-response}{}
\addcontentsline{toc}{section}{UC-18 --- Mark True Response}
\section*{UC-18 --- Mark True Response}
\par\medskip

\textbf{Version:} \texttt{1.0}
\par\medskip

\textbf{Performed by:} [Nguyen Le Quoc Huy] | \textbf{Edited by:} [Nguyen Le Quoc Huy]
\par\medskip

\begingroup
\small
\renewcommand{\baselinestretch}{1.0}\selectfont
\renewcommand{\arraystretch}{1.15}
\setlength{\tabcolsep}{6pt}
\begin{longtable}{|>{\raggedright\arraybackslash}p{0.10\textwidth}|>{\raggedright\arraybackslash}p{0.83\textwidth}|}
\hline
\textbf{Field} & \textbf{Value} \\ \hline
\endhead
Use Case ID & UC-18 \\ \hline
Name & Mark True Response \\ \hline
Primary Actor & Post Author \\ \hline
Supporting Actors & Answer Author \\ \hline
Feature Group & Forum Interaction \\ \hline
Priority & Not explicitly stated \\ \hline
Related User Stories & US-20 \\ \hline
Related Requirements & FR-13.2, FR-17.3, FR-17.4, FR-18.2, FR-20.1 to FR-20.5, FR-38.1; NFR-PERF-3, NFR-PERF-4, NFR-PERF-6, NFR-SEC-1, NFR-SEC-6, NFR-PRIV-7 \\ \hline
Source Documents & \texttt{proposals/userstory.md} (section 6.3); \texttt{proposals/functional\_requirement.md} (US-20); \texttt{proposals/list\_post.md}; \texttt{proposals/post\_management.md}; \texttt{proposals/API.md}; \texttt{proposals/response\_code.md}; \texttt{proposals/non\_functional\_requirement.md}; \texttt{proposals/feature\_flow.md}; \texttt{proposals/screen\_flow.md} \\ \hline
\end{longtable}
\endgroup
\par\medskip

\hypertarget{sec:part-a-use-case-model}{}
\addcontentsline{toc}{subsection}{Part A --- Use Case Model}
\subsection*{Part A --- Use Case Model}
\par\medskip

\hypertarget{sec:a1-focused-use-case-diagram}{}
\addcontentsline{toc}{subsubsection}{A1. Focused Use Case Diagram}
\subsubsection*{A1. Focused Use Case Diagram}
\par\medskip

\begin{center}
\includegraphics[width=0.85\textwidth,height=0.75\textheight,keepaspectratio]{"figures/UC-18-mark-true-response_1.png"}
\end{center}
\par\medskip

\textbf{Diagram notation.} Person shape = human actor; cloud = supporting system; blue = selected use case; solid line = association; labelled solid arrow = generalization; dashed arrow = include or extend.
\par\medskip

\hypertarget{sec:a2-actor-register}{}
\addcontentsline{toc}{subsubsection}{A2. Actor Register}
\subsubsection*{A2. Actor Register}
\par\medskip

\begingroup
\small
\renewcommand{\baselinestretch}{1.0}\selectfont
\renewcommand{\arraystretch}{1.15}
\setlength{\tabcolsep}{6pt}
\begin{longtable}{|>{\raggedright\arraybackslash}p{0.10\textwidth}|>{\raggedright\arraybackslash}p{0.17\textwidth}|>{\raggedright\arraybackslash}p{0.65\textwidth}|}
\hline
\textbf{Actor} & \textbf{Type} & \textbf{Description} \\ \hline
\endhead
Post Author & Primary human actor & The author of a question post who marks one comment or answer as the true response. \\ \hline
Answer Author & Secondary human actor & The author of the selected answer who is notified that the answer was accepted. \\ \hline
\end{longtable}
\endgroup
\par\medskip

\hypertarget{sec:a3-relationships}{}
\addcontentsline{toc}{subsubsection}{A3. Relationships}
\subsubsection*{A3. Relationships}
\par\medskip

\hypertarget{sec:a3-1-associations}{}
\paragraph*{A3.1 Associations}
\par\medskip

\begingroup
\small
\renewcommand{\baselinestretch}{1.0}\selectfont
\renewcommand{\arraystretch}{1.15}
\setlength{\tabcolsep}{6pt}
\begin{longtable}{|>{\raggedright\arraybackslash}p{0.12\textwidth}|>{\raggedright\arraybackslash}p{0.22\textwidth}|>{\raggedright\arraybackslash}p{0.10\textwidth}|>{\raggedright\arraybackslash}p{0.48\textwidth}|}
\hline
\textbf{Actor} & \textbf{Use Case} & \textbf{Type} & \textbf{Description} \\ \hline
\endhead
Post Author & UC-18 Mark True Response & Primary & Selects or changes the accepted comment or answer. \\ \hline
Answer Author & UC-18 Mark True Response & Secondary & Receives the acceptance notification after selection. \\ \hline
\end{longtable}
\endgroup
\par\medskip

\hypertarget{sec:a3-2-generalization-and-specialization}{}
\paragraph*{A3.2 Generalization and Specialization}
\par\medskip

\begingroup
\small
\renewcommand{\baselinestretch}{1.0}\selectfont
\renewcommand{\arraystretch}{1.15}
\setlength{\tabcolsep}{6pt}
\begin{longtable}{|>{\raggedright\arraybackslash}p{0.09\textwidth}|>{\raggedright\arraybackslash}p{0.09\textwidth}|>{\raggedright\arraybackslash}p{0.09\textwidth}|>{\raggedright\arraybackslash}p{0.65\textwidth}|}
\hline
\textbf{Child} & \textbf{Parent} & \textbf{Type} & \textbf{Inherited Behavior or Association} \\ \hline
\endhead
None & None & None & No actor generalization or specialization is named for US-20 in the specified proposals. \\ \hline
\end{longtable}
\endgroup
\par\medskip

\hypertarget{sec:a3-3-include}{}
\paragraph*{A3.3 Include}
\par\medskip

\begingroup
\small
\renewcommand{\baselinestretch}{1.0}\selectfont
\renewcommand{\arraystretch}{1.15}
\setlength{\tabcolsep}{6pt}
\begin{longtable}{|>{\raggedright\arraybackslash}p{0.10\textwidth}|>{\raggedright\arraybackslash}p{0.11\textwidth}|>{\raggedright\arraybackslash}p{0.71\textwidth}|}
\hline
\textbf{Base Use Case} & \textbf{Included Use Case} & \textbf{Meaning} \\ \hline
\endhead
None & None & No distinct included use case is named for US-20 in the specified proposals; no new I-xx ID is introduced. \\ \hline
\end{longtable}
\endgroup
\par\medskip

\hypertarget{sec:a3-4-extend}{}
\paragraph*{A3.4 Extend}
\par\medskip

\begingroup
\small
\renewcommand{\baselinestretch}{1.0}\selectfont
\renewcommand{\arraystretch}{1.15}
\setlength{\tabcolsep}{6pt}
\begin{longtable}{|>{\raggedright\arraybackslash}p{0.10\textwidth}|>{\raggedright\arraybackslash}p{0.09\textwidth}|>{\raggedright\arraybackslash}p{0.09\textwidth}|>{\raggedright\arraybackslash}p{0.62\textwidth}|}
\hline
\textbf{Extending Use Case} & \textbf{Base Use Case} & \textbf{Condition} & \textbf{Added Behavior} \\ \hline
\endhead
None & None & None & No distinct extending use case is named for US-20 in the specified proposals; no new E-xx ID is introduced. \\ \hline
\end{longtable}
\endgroup
\par\medskip

\hypertarget{sec:a4-model-notes}{}
\addcontentsline{toc}{subsubsection}{A4. Model Notes}
\subsubsection*{A4. Model Notes}
\par\medskip

\begin{enumerate}
\item UC-18 applies to a question post; the post author is the only actor allowed to select the true response.
\item The Answer Author association models the notification recipient named by \texttt{AC-20.5} and \texttt{FR-20.5}; no separate notification service actor is introduced from unrelated sources.
\item Marking and changing are the same use case because \texttt{AC-20.4} and \texttt{FR-20.4} describe a change to the accepted response.
\item The marked response is a comment or answer under \texttt{FR-20.1}; \texttt{API.md} identifies the target with \texttt{comment\_id}. Priority and a dedicated true-response UI route are not specified in proposals.
\end{enumerate}
\par\medskip

\hypertarget{sec:part-b-use-case-specification}{}
\addcontentsline{toc}{subsection}{Part B --- Use Case Specification}
\subsection*{Part B --- Use Case Specification}
\par\medskip

\hypertarget{sec:b1-brief-description}{}
\addcontentsline{toc}{subsubsection}{B1. Brief Description}
\subsubsection*{B1. Brief Description}
\par\medskip

A post author marks one comment or answer as the true response so that other students can quickly find the accepted solution. The post is marked as solved, and the true response is displayed prominently under the post.
\par\medskip

The author can change the true response when necessary. The author of the selected answer is notified that the answer was accepted.
\par\medskip

\hypertarget{sec:b2-preconditions}{}
\addcontentsline{toc}{subsubsection}{B2. Preconditions}
\subsubsection*{B2. Preconditions}
\par\medskip

\begin{enumerate}
\item The post author is authenticated with an active session.
\item The post is a question post owned by the acting author.
\item At least one comment or answer exists on the post.
\item The selected comment or answer is visible and available for acceptance.
\end{enumerate}
\par\medskip

\hypertarget{sec:b3-postconditions}{}
\addcontentsline{toc}{subsubsection}{B3. Postconditions}
\subsubsection*{B3. Postconditions}
\par\medskip

\begin{enumerate}
\item \textbf{Success:} The selected comment or answer is recorded as the post's true response.
\item \textbf{Success:} The post is marked as solved and the true response is displayed prominently under the post.
\item \textbf{Success:} The author of the selected answer is notified that the answer was accepted.
\item \textbf{Failure:} The true response and solved state remain unchanged, and the post author receives a clear error.
\end{enumerate}
\par\medskip

\hypertarget{sec:b4-main-success-scenario-basic-flow}{}
\addcontentsline{toc}{subsubsection}{B4. Main Success Scenario (Basic Flow)}
\subsubsection*{B4. Main Success Scenario (Basic Flow)}
\par\medskip

\begin{enumerate}
\item The post author opens the question post detail page.
\item The system displays the post, its comments and answers, and the current true response when present.
\item The system displays \textbf{Mark as true response} only to the post author.
\item The post author selects a comment or answer.
\item The system sends \texttt{comment\_id} to the true-response endpoint.
\item The backend verifies authentication, post ownership, post existence, and comment existence.
\item The system records the selected response and marks the post solved.
\item The system returns \texttt{solved: true}.
\item The system displays the true response prominently under the post and updates the solved status.
\item The system notifies the author of the selected answer that the answer was accepted.
\end{enumerate}
\par\medskip

\hypertarget{sec:b5-alternative-flows}{}
\addcontentsline{toc}{subsubsection}{B5. Alternative Flows}
\subsubsection*{B5. Alternative Flows}
\par\medskip

The story has no distinct alternative beyond source-backed variants of the same action.
\par\medskip

\hypertarget{sec:a1-change-the-true-response}{}
\paragraph*{A1. Change the true response}
\par\medskip

The post author selects a different comment or answer. The system replaces the previous true response with the new selection, keeps the post solved, and notifies the newly selected answer author.
\par\medskip

\hypertarget{sec:a2-mark-a-comment-instead-of-an-answer}{}
\paragraph*{A2. Mark a comment instead of an answer}
\par\medskip

At basic flow step 4, the author selects a comment rather than an answer. \texttt{FR-20.1} allows a comment or answer to be marked, and the API target remains a \texttt{comment\_id}.
\par\medskip

\hypertarget{sec:b6-exception-flows}{}
\addcontentsline{toc}{subsubsection}{B6. Exception Flows}
\subsubsection*{B6. Exception Flows}
\par\medskip

\hypertarget{sec:e1-post-or-comment-not-found}{}
\paragraph*{E1. Post or comment not found}
\par\medskip

The system returns \texttt{404} with \texttt{POST\_NOT\_FOUND} (\texttt{48001}) or \texttt{COMMENT\_NOT\_FOUND} (\texttt{48002}) and refreshes the available post discussion.
\par\medskip

\hypertarget{sec:e2-caller-is-not-the-post-owner}{}
\paragraph*{E2. Caller is not the post owner}
\par\medskip

The system returns \texttt{403} with \texttt{NOT\_POST\_OWNER} (\texttt{48007}) or \texttt{FORBIDDEN} (\texttt{48010}) and does not change the true response.
\par\medskip

\hypertarget{sec:e3-unauthorized-request}{}
\paragraph*{E3. Unauthorized request}
\par\medskip

The system returns \texttt{401} and requires authentication before the true response is changed.
\par\medskip

\hypertarget{sec:e4-already-true-response-conflict}{}
\paragraph*{E4. Already true response conflict}
\par\medskip

\texttt{ALREADY\_TRUE\_RESPONSE} (\texttt{48008}) exists in \texttt{response\_code.md}, but it is not listed in the \texttt{PATCH /forum/posts/:post\_id/true-response} \texttt{403} contract. No source-defined status mapping is available for this case.
\par\medskip

\hypertarget{sec:e5-notification-delivery-failure}{}
\paragraph*{E5. Notification delivery failure}
\par\medskip

The true response remains marked and the post remains solved. Notification retry or recording behavior is not specified in proposals.
\par\medskip

\hypertarget{sec:e6-server-or-database-failure}{}
\paragraph*{E6. Server or database failure}
\par\medskip

The system returns \texttt{500} with \texttt{DATABASE\_ERROR} or \texttt{UNKNOWN\_ERROR} (\texttt{45000}), leaves the previous true-response state authoritative, and offers retry where supported.
\par\medskip

\hypertarget{sec:b7-business-rules}{}
\addcontentsline{toc}{subsubsection}{B7. Business Rules}
\subsubsection*{B7. Business Rules}
\par\medskip

\begin{enumerate}
\item \texttt{AC-20.1} / \texttt{FR-20.1}: The post author can mark one comment or answer as the true response.
\item \texttt{AC-20.2} / \texttt{FR-20.2}: The post is marked as solved when a true response is selected.
\item \texttt{AC-20.3} / \texttt{FR-20.3}: The true response is displayed prominently under the post.
\item \texttt{AC-20.4} / \texttt{FR-20.4}: The author can change the true response when necessary.
\item \texttt{AC-20.5} / \texttt{FR-20.5}: The author of the selected answer is notified that the answer was accepted.
\item Only the post author may perform the action; the backend enforces ownership under \texttt{FR-20.1} and the \texttt{NOT\_POST\_OWNER} (\texttt{48007}) contract.
\item The true-response target is identified by \texttt{comment\_id: string}.
\item Privacy settings apply consistently to the post and answer display (\texttt{NFR-PRIV-7}).
\end{enumerate}
\par\medskip

\hypertarget{sec:b8-special-requirements}{}
\addcontentsline{toc}{subsubsection}{B8. Special Requirements}
\subsubsection*{B8. Special Requirements}
\par\medskip

\begingroup
\small
\renewcommand{\baselinestretch}{1.0}\selectfont
\renewcommand{\arraystretch}{1.15}
\setlength{\tabcolsep}{6pt}
\begin{longtable}{|>{\raggedright\arraybackslash}p{0.14\textwidth}|>{\raggedright\arraybackslash}p{0.79\textwidth}|}
\hline
\textbf{ID} & \textbf{Requirement} \\ \hline
\endhead
NFR-PERF-3, NFR-PERF-4 & Post detail and comments display within 2 seconds, and a true-response mark is reflected in the forum view within 2 seconds. \\ \hline
NFR-PERF-6 & Forum notifications, including true-response marks, must be delivered within 2 seconds under normal conditions. \\ \hline
NFR-SEC-1 & Authentication is required before accessing protected interactions. \\ \hline
NFR-SEC-6 & Server-side authorization checks must prevent unauthorized true-response changes. \\ \hline
NFR-PRIV-7 & Privacy settings apply consistently to forum posts, comments, member lists, and search results. \\ \hline
Design constraint & The acceptance notification must not roll back a successful true-response mark when delivery fails. \\ \hline
\end{longtable}
\endgroup
\par\medskip

\hypertarget{sec:b9-extension-points}{}
\addcontentsline{toc}{subsubsection}{B9. Extension Points}
\subsubsection*{B9. Extension Points}
\par\medskip

\begingroup
\small
\renewcommand{\baselinestretch}{1.0}\selectfont
\renewcommand{\arraystretch}{1.15}
\setlength{\tabcolsep}{6pt}
\begin{longtable}{|>{\raggedright\arraybackslash}p{0.37\textwidth}|>{\raggedright\arraybackslash}p{0.56\textwidth}|}
\hline
\textbf{Extension Point} & \textbf{Location in Flow of Events} \\ \hline
\endhead
Author-only control visibility & Basic flow step 3 \\ \hline
Select response & Basic flow step 4; alternatives A1-A2 \\ \hline
Ownership check & Basic flow step 6; exception E2 \\ \hline
Solved-state update & Basic flow steps 7-9 \\ \hline
Acceptance notification & Basic flow step 10; exception E5 \\ \hline
Failure recovery & Basic flow steps 5-7; exceptions E1-E4 and E6 \\ \hline
\end{longtable}
\endgroup
\par\medskip

\hypertarget{sec:b10-interface-and-data-notes}{}
\addcontentsline{toc}{subsubsection}{B10. Interface and Data Notes}
\subsubsection*{B10. Interface and Data Notes}
\par\medskip

\begin{itemize}
\item \textbf{Screen/flow path:} \texttt{screen\_flow.md} gives \texttt{Forum -$>$ University -$>$ Major -$>$ Course -$>$ Post}, with \texttt{Trending -$>$ Post}. \texttt{feature\_flow.md} gives \texttt{Post -$>$ Comment -$>$ True Response}, \texttt{List Post -$>$ Solved / Good -$>$ Comment / Answer}, and \texttt{List Post -$>$ Post -$>$ Comment / Solved / Vote / Save / Report / Moderate}.
\item \textbf{UI components:} \texttt{list\_post.md} and \texttt{post\_management.md} recommend \texttt{MarkTrueResponseButton} with label \textbf{Mark as true response}, \texttt{TrueResponsePanel}, \texttt{PostSolvedBadge} with \texttt{Solved} / \texttt{Unsolved}, and \texttt{PermissionBasedActionBar}.
\item \textbf{Detail load:} \texttt{GET /forum/posts/:post\_id} returns \texttt{post}, \texttt{comments}, \texttt{true\_response: object | null}, \texttt{user\_vote}, and \texttt{saved} with \texttt{200}; failures are \texttt{404 POST\_NOT\_FOUND} (\texttt{48001}), \texttt{401}, and \texttt{500 (DATABASE\_ERROR | UNKNOWN\_ERROR)}. Database function \texttt{studitopia.get\_forum\_post} also lists \texttt{FORBIDDEN: 48010}.
\item \textbf{API endpoint and failures:} \texttt{PATCH /forum/posts/:post\_id/true-response} with body \texttt{\{ comment\_id: string \}} returns \texttt{200} with \texttt{data: \{ solved: boolean \}}; failures are \texttt{404} with \texttt{(POST\_NOT\_FOUND | COMMENT\_NOT\_FOUND)[]} (\texttt{48001}, \texttt{48002}), \texttt{403} with \texttt{(NOT\_POST\_OWNER | FORBIDDEN)[]} (\texttt{48007}, \texttt{48010}), \texttt{401}, and \texttt{500} with \texttt{(DATABASE\_ERROR | UNKNOWN\_ERROR)[]}.
\item \textbf{Database contract and response-code gap:} \texttt{studitopia.mark\_true\_response(puser\_id: uuid, ppost\_id: uuid, pcomment\_id: uuid)} returns \texttt{\{solved: boolean\}} and declares \texttt{UNKNOWN\_ERROR: 45000}, \texttt{POST\_NOT\_FOUND: 48001}, \texttt{COMMENT\_NOT\_FOUND: 48002}, \texttt{NOT\_POST\_OWNER: 48007}, and \texttt{FORBIDDEN: 48010}. \texttt{response\_code.md} defines \texttt{48008 ALREADY\_TRUE\_RESPONSE} in the forum, database, and backend registers, but \texttt{API.md} does not list it in the true-response endpoint contract.
\item \textbf{List display and notifications:} \texttt{GET /forum/posts} returns \texttt{solved: boolean} in each post summary so solved status is visible while browsing. \texttt{FR-38.1} and \texttt{NFR-PERF-6} include true-response marks in forum notifications; no notification endpoint or payload is defined in \texttt{API.md}.
\end{itemize}
\par\medskip

\hypertarget{sec:b11-source-alignment-notes}{}
\addcontentsline{toc}{subsubsection}{B11. Source Alignment Notes}
\subsubsection*{B11. Source Alignment Notes}
\par\medskip

\begin{itemize}
\item UC-04 covers post publication and lists true-response eligibility as downstream context; UC-18 covers marking an existing response and does not repeat discussion-post creation. UC-05 covers comment and answer creation and notification of the post author; UC-18 begins after those responses exist.
\item UC-05 already notes the same \texttt{ALREADY\_TRUE\_RESPONSE} (\texttt{48008}) gap: the code exists in \texttt{response\_code.md} but is missing from the \texttt{PATCH /forum/posts/:post\_id/true-response} \texttt{403} contract. This document preserves that gap rather than assigning a status code.
\item \texttt{FR-20.1} allows a comment or answer to be marked, while the API identifies the target only as \texttt{comment\_id}; no separate answer identifier is specified. \texttt{AC-20.4} permits changing the true response, but no distinct change endpoint or dedicated code is defined; the same PATCH endpoint is the only source-backed contract.
\item Notification delivery is required by \texttt{FR-20.5}, while user notification preferences are a cross-cutting concern; no source specifies how mandatory acceptance notifications interact with opt-outs. Route names and component names in \texttt{list\_post.md} and \texttt{post\_management.md} are recommendations; \texttt{API.md} defines the exact endpoint and codes used above. No distinct included, extending, or specialized use case is named for US-20, so no new I-xx/E-xx/G-xx ID is introduced.
\end{itemize}
\par\medskip

\par\medskip
\hypertarget{sec:ui-prototype-specimens}{}
\addcontentsline{toc}{subsection}{UI Prototype Specimens}
\subsection*{UI Prototype Specimens}
\par\medskip
\begin{figure}[H]
\centering
\includegraphics[width=0.75\textwidth,height=0.42\textheight,keepaspectratio]{"../../prototype/dijkstra_question_desktop.png"}
\par\vspace{2pt}
{\footnotesize \textbf{UI Prototype:} Dijkstra question - Desktop}
\end{figure}
\par\smallskip

\pagebreak
\hypertarget{sec:uc-19-save-post}{}
\addcontentsline{toc}{section}{UC-19 --- Save Post}
\section*{UC-19 --- Save Post}
\par\medskip

\textbf{Version:} \texttt{1.0}
\par\medskip

\textbf{Performed by:} [Nguyen Le Quoc Huy] | \textbf{Edited by:} [Nguyen Le Quoc Huy]
\par\medskip

\begingroup
\small
\renewcommand{\baselinestretch}{1.0}\selectfont
\renewcommand{\arraystretch}{1.15}
\setlength{\tabcolsep}{6pt}
\begin{longtable}{|>{\raggedright\arraybackslash}p{0.10\textwidth}|>{\raggedright\arraybackslash}p{0.83\textwidth}|}
\hline
\textbf{Field} & \textbf{Value} \\ \hline
\endhead
Use Case ID & UC-19 \\ \hline
Name & Save Post \\ \hline
Primary Actor & University Student \\ \hline
Supporting Actors & None \\ \hline
Feature Group & Forum Interaction \\ \hline
Priority & Not specified in proposals \\ \hline
Related User Stories & US-21 \\ \hline
Related Requirements & FR-21.1, FR-21.2, FR-21.3, FR-21.4; AC-21.1, AC-21.2, AC-21.3, AC-21.4; NFR-PERF-1, NFR-SEC-1, NFR-SEC-6, NFR-PRIV-7 \\ \hline
Source Documents & \texttt{proposals/userstory.md} §6.4; \texttt{proposals/functional\_requirement.md} US-21; \texttt{proposals/post\_management.md}; \texttt{proposals/list\_post.md}; \texttt{proposals/API.md}; \texttt{proposals/response\_code.md}; \texttt{proposals/non\_functional\_requirement.md}; \texttt{proposals/settings.md}; \texttt{proposals/feature\_flow.md}; \texttt{proposals/screen\_flow.md} \\ \hline
\end{longtable}
\endgroup
\par\medskip

\hypertarget{sec:part-a-use-case-model}{}
\addcontentsline{toc}{subsection}{Part A --- Use Case Model}
\subsection*{Part A --- Use Case Model}
\par\medskip

\hypertarget{sec:a1-focused-use-case-diagram}{}
\addcontentsline{toc}{subsubsection}{A1. Focused Use Case Diagram}
\subsubsection*{A1. Focused Use Case Diagram}
\par\medskip

\begin{center}
\includegraphics[width=0.85\textwidth,height=0.75\textheight,keepaspectratio]{"figures/UC-19-save-post_1.png"}
\end{center}
\par\medskip

\textbf{Diagram notation.} Person shape = human actor; cloud = supporting system; blue = selected use case; solid line = association; labelled solid arrow = generalization; dashed arrow = include or extend.
\par\medskip

\hypertarget{sec:a2-actor-register}{}
\addcontentsline{toc}{subsubsection}{A2. Actor Register}
\subsubsection*{A2. Actor Register}
\par\medskip

\begingroup
\small
\renewcommand{\baselinestretch}{1.0}\selectfont
\renewcommand{\arraystretch}{1.15}
\setlength{\tabcolsep}{6pt}
\begin{longtable}{|>{\raggedright\arraybackslash}p{0.11\textwidth}|>{\raggedright\arraybackslash}p{0.11\textwidth}|>{\raggedright\arraybackslash}p{0.70\textwidth}|}
\hline
\textbf{Actor} & \textbf{Type} & \textbf{Description} \\ \hline
\endhead
University Student & Primary human actor & A registered student who saves a useful discussion or material post and returns to it later through the saved list. \\ \hline
\end{longtable}
\endgroup
\par\medskip

\hypertarget{sec:a3-relationships}{}
\addcontentsline{toc}{subsubsection}{A3. Relationships}
\subsubsection*{A3. Relationships}
\par\medskip

\hypertarget{sec:a3-1-associations}{}
\paragraph*{A3.1 Associations}
\par\medskip

\begingroup
\small
\renewcommand{\baselinestretch}{1.0}\selectfont
\renewcommand{\arraystretch}{1.15}
\setlength{\tabcolsep}{6pt}
\begin{longtable}{|>{\raggedright\arraybackslash}p{0.15\textwidth}|>{\raggedright\arraybackslash}p{0.13\textwidth}|>{\raggedright\arraybackslash}p{0.10\textwidth}|>{\raggedright\arraybackslash}p{0.53\textwidth}|}
\hline
\textbf{Actor} & \textbf{Use Case} & \textbf{Type} & \textbf{Description} \\ \hline
\endhead
University Student & UC-19 Save Post & Primary & Saves or removes a post and views the user-scoped saved list. \\ \hline
\end{longtable}
\endgroup
\par\medskip

\hypertarget{sec:a3-2-generalization-and-specialization}{}
\paragraph*{A3.2 Generalization and Specialization}
\par\medskip

\begingroup
\small
\renewcommand{\baselinestretch}{1.0}\selectfont
\renewcommand{\arraystretch}{1.15}
\setlength{\tabcolsep}{6pt}
\begin{longtable}{|>{\raggedright\arraybackslash}p{0.09\textwidth}|>{\raggedright\arraybackslash}p{0.09\textwidth}|>{\raggedright\arraybackslash}p{0.09\textwidth}|>{\raggedright\arraybackslash}p{0.65\textwidth}|}
\hline
\textbf{Child} & \textbf{Parent} & \textbf{Type} & \textbf{Inherited Behavior or Association} \\ \hline
\endhead
None & None & None & No actor generalization or specialization is named in the proposal sources for this story. \\ \hline
\end{longtable}
\endgroup
\par\medskip

\hypertarget{sec:a3-3-include}{}
\paragraph*{A3.3 Include}
\par\medskip

\begingroup
\small
\renewcommand{\baselinestretch}{1.0}\selectfont
\renewcommand{\arraystretch}{1.15}
\setlength{\tabcolsep}{6pt}
\begin{longtable}{|>{\raggedright\arraybackslash}p{0.11\textwidth}|>{\raggedright\arraybackslash}p{0.14\textwidth}|>{\raggedright\arraybackslash}p{0.67\textwidth}|}
\hline
\textbf{Base Use Case} & \textbf{Included Use Case} & \textbf{Meaning} \\ \hline
\endhead
None & None & No distinct included use case is named in the proposal sources for this story. \\ \hline
\end{longtable}
\endgroup
\par\medskip

\hypertarget{sec:a3-4-extend}{}
\paragraph*{A3.4 Extend}
\par\medskip

\begingroup
\small
\renewcommand{\baselinestretch}{1.0}\selectfont
\renewcommand{\arraystretch}{1.15}
\setlength{\tabcolsep}{6pt}
\begin{longtable}{|>{\raggedright\arraybackslash}p{0.14\textwidth}|>{\raggedright\arraybackslash}p{0.10\textwidth}|>{\raggedright\arraybackslash}p{0.10\textwidth}|>{\raggedright\arraybackslash}p{0.58\textwidth}|}
\hline
\textbf{Extending Use Case} & \textbf{Base Use Case} & \textbf{Condition} & \textbf{Added Behavior} \\ \hline
\endhead
None & None & None & No distinct extending use case is named in the proposal sources for this story. \\ \hline
\end{longtable}
\endgroup
\par\medskip

\hypertarget{sec:a4-model-notes}{}
\addcontentsline{toc}{subsubsection}{A4. Model Notes}
\subsubsection*{A4. Model Notes}
\par\medskip

\begin{enumerate}
\item In scope: saving a post, removing a saved post, listing saved posts, and keeping saved posts available across forum context changes.
\item Out of scope: voting, commenting, true-response marking, reporting, and post creation, which are separate forum requirements.
\item No new I-xx, E-xx, or G-xx identifier is created; the proposal sources name no distinct included, extending, or specialized use case for saving posts.
\item \texttt{API.md} defines a save/unsave endpoint and a \texttt{saved} field on post detail, but no dedicated saved-list endpoint; this gap is retained rather than filled with an invented route.
\item \texttt{SavedPostsPanel}, \texttt{SavedPostsLink}, \texttt{SavedPostList}, and suggested routes such as \texttt{/forum/saved} are source-backed UI recommendations, not confirmed API routes or FR/AC fields.
\end{enumerate}
\par\medskip

\hypertarget{sec:part-b-use-case-specification}{}
\addcontentsline{toc}{subsection}{Part B --- Use Case Specification}
\subsection*{Part B --- Use Case Specification}
\par\medskip

\hypertarget{sec:b1-brief-description}{}
\addcontentsline{toc}{subsubsection}{B1. Brief Description}
\subsubsection*{B1. Brief Description}
\par\medskip

A university student saves a post so that a useful discussion or material can be found again later. The saved list is user-scoped and remains available when the student changes forum context.
\par\medskip

The trigger is selecting the saved-state control on a visible post. The successful outcome is an updated \texttt{saved} state, a saved post added to or removed from the user's list, and continued access to saved posts across university, major, and course forum contexts.
\par\medskip

\hypertarget{sec:b2-preconditions}{}
\addcontentsline{toc}{subsubsection}{B2. Preconditions}
\subsubsection*{B2. Preconditions}
\par\medskip

\begin{enumerate}
\item The student is authenticated with an active session.
\item The student can view the selected post.
\item The post exists and has not been deleted or made unavailable.
\item The system can determine the current user's saved state for the post.
\end{enumerate}
\par\medskip

\hypertarget{sec:b3-postconditions}{}
\addcontentsline{toc}{subsubsection}{B3. Postconditions}
\subsubsection*{B3. Postconditions}
\par\medskip

\begin{enumerate}
\item \textbf{Success:} The post is saved for the requesting user and the saved state is returned as \texttt{saved: boolean}.
\item \textbf{Success:} The system provides the user's saved-post list and keeps saved posts available regardless of current forum context.
\item \textbf{Success:} A removal leaves the post unsaved and no longer listed as saved for that user.
\item \textbf{Failure:} The saved state is unchanged and the student receives a clear error.
\end{enumerate}
\par\medskip

\hypertarget{sec:b4-main-success-scenario-basic-flow}{}
\addcontentsline{toc}{subsubsection}{B4. Main Success Scenario (Basic Flow)}
\subsubsection*{B4. Main Success Scenario (Basic Flow)}
\par\medskip

\begin{enumerate}
\item The student opens a forum post list or post detail view.
\item The system displays the post and the actions permitted for the student.
\item The system shows the current saved state for the post.
\item The student selects the save control.
\item The frontend sends the save request to the backend for the current user and target post.
\item The backend verifies the session, post visibility, and authorization, then sets the saved state for that user.
\item The backend returns \texttt{saved: boolean}.
\item The UI updates the saved state and the saved-post list when that list is displayed.
\end{enumerate}
\par\medskip

\hypertarget{sec:b5-alternative-flows}{}
\addcontentsline{toc}{subsubsection}{B5. Alternative Flows}
\subsubsection*{B5. Alternative Flows}
\par\medskip

\hypertarget{sec:a1-remove-a-saved-post}{}
\paragraph*{A1. Remove a saved post}
\par\medskip

At basic flow step 4, the student selects the saved-state control for an already saved post. The system unsaves the post for that user and returns the updated saved state.
\par\medskip

\hypertarget{sec:a2-change-forum-context}{}
\paragraph*{A2. Change forum context}
\par\medskip

The student changes university, major, or course forum context and later opens the saved list. The saved posts remain available because saved state is user-scoped and persists across context changes.
\par\medskip

\hypertarget{sec:a3-empty-saved-list}{}
\paragraph*{A3. Empty saved list}
\par\medskip

The student opens \textbf{Saved Posts} when no posts are saved. The system shows the source-recommended \texttt{EmptySavedPostsState} with "No saved posts" and a browse link.
\par\medskip

\hypertarget{sec:a4-open-a-saved-post}{}
\paragraph*{A4. Open a saved post}
\par\medskip

From the saved list, the student opens a saved post. The system displays the post detail if the post is still visible and available.
\par\medskip

\hypertarget{sec:b6-exception-flows}{}
\addcontentsline{toc}{subsubsection}{B6. Exception Flows}
\subsubsection*{B6. Exception Flows}
\par\medskip

\hypertarget{sec:e1-post-is-unavailable}{}
\paragraph*{E1. Post is unavailable}
\par\medskip

If the target post does not exist, the backend reports \texttt{POST\_NOT\_FOUND} (\texttt{48001}) under \texttt{404} in \texttt{API.md}. The UI leaves the saved state unchanged and shows a post-unavailable state.
\par\medskip

\hypertarget{sec:e2-student-is-unauthenticated}{}
\paragraph*{E2. Student is unauthenticated}
\par\medskip

If the session is missing or invalid, the backend reports \texttt{UNAUTHORIZED} under \texttt{401}. The system denies the save or remove action and requires authentication.
\par\medskip

\hypertarget{sec:e3-authorization-failure}{}
\paragraph*{E3. Authorization failure}
\par\medskip

\texttt{studitopia.save\_forum\_post(puser\_id, ppost\_id)} lists \texttt{FORBIDDEN} (\texttt{48010}), and \texttt{response\_code.md} defines \texttt{48010 FORBIDDEN}. \texttt{API.md} does not specify an HTTP status for that code on this endpoint, so no status is invented.
\par\medskip

\hypertarget{sec:e4-server-or-database-failure}{}
\paragraph*{E4. Server or database failure}
\par\medskip

If saving fails, the backend reports \texttt{INTERNAL\_SERVER\_ERROR} under \texttt{500} with \texttt{(DATABASE\_ERROR | UNKNOWN\_ERROR)[]}. The UI preserves the prior saved state and offers retry where supported.
\par\medskip

\hypertarget{sec:b7-business-rules}{}
\addcontentsline{toc}{subsubsection}{B7. Business Rules}
\subsubsection*{B7. Business Rules}
\par\medskip

\begin{enumerate}
\item Users can save a post under FR-21.1 and AC-21.1.
\item The system provides a list of the user's saved posts under FR-21.2 and AC-21.2.
\item Users can remove a post from their saved list under FR-21.3 and AC-21.3.
\item Saved posts remain available regardless of current forum context under FR-21.4 and AC-21.4.
\item Save state is user-scoped; \texttt{studitopia.save\_forum\_post} saves or unsaves a post for the requesting user.
\item Only posts visible to the requesting user are returned by \texttt{studitopia.list\_forum\_posts}, and \texttt{studitopia.get\_forum\_post} respects privacy and moderation visibility rules.
\item No numeric limit on saved posts is specified in proposals.
\end{enumerate}
\par\medskip

\hypertarget{sec:b8-special-requirements}{}
\addcontentsline{toc}{subsubsection}{B8. Special Requirements}
\subsubsection*{B8. Special Requirements}
\par\medskip

\begingroup
\small
\renewcommand{\baselinestretch}{1.0}\selectfont
\renewcommand{\arraystretch}{1.15}
\setlength{\tabcolsep}{6pt}
\begin{longtable}{|>{\raggedright\arraybackslash}p{0.14\textwidth}|>{\raggedright\arraybackslash}p{0.79\textwidth}|}
\hline
\textbf{ID} & \textbf{Requirement} \\ \hline
\endhead
NFR-PERF-1 & Standard user actions must respond within 2 seconds. \\ \hline
NFR-SEC-1 & Authentication is required before accessing protected user resources. \\ \hline
NFR-SEC-6 & Server-side authorization checks must prevent unauthorized access. \\ \hline
NFR-PRIV-7 & Privacy settings must apply consistently to forum posts, comments, member lists, and search results. \\ \hline
Design constraint & Saved state must not be cleared when the student switches forum context. \\ \hline
\end{longtable}
\endgroup
\par\medskip

\hypertarget{sec:b9-extension-points}{}
\addcontentsline{toc}{subsubsection}{B9. Extension Points}
\subsubsection*{B9. Extension Points}
\par\medskip

\begingroup
\small
\renewcommand{\baselinestretch}{1.0}\selectfont
\renewcommand{\arraystretch}{1.15}
\setlength{\tabcolsep}{6pt}
\begin{longtable}{|>{\raggedright\arraybackslash}p{0.39\textwidth}|>{\raggedright\arraybackslash}p{0.54\textwidth}|}
\hline
\textbf{Extension Point} & \textbf{Location in Flow of Events} \\ \hline
\endhead
Save or unsave decision & Basic flow steps 3-5; alternatives A1 and A4 \\ \hline
Saved-list display & Basic flow step 8; alternatives A2-A3 \\ \hline
Unavailable post handling & Basic flow step 2; exception E1 \\ \hline
Authentication and authorization & Basic flow step 6; exceptions E2-E3 \\ \hline
Backend failure recovery & Basic flow step 7; exception E4 \\ \hline
\end{longtable}
\endgroup
\par\medskip

\hypertarget{sec:b10-interface-and-data-notes}{}
\addcontentsline{toc}{subsubsection}{B10. Interface and Data Notes}
\subsubsection*{B10. Interface and Data Notes}
\par\medskip

\begin{itemize}
\item \textbf{Confirmed requirement:} \texttt{feature\_flow.md} and \texttt{userstory.md} place \textbf{Save} after \texttt{Post}; \texttt{list\_post.md} places \texttt{Post → Save} in the post-action path.
\item \textbf{Recommendation UI:} \texttt{post\_management.md} proposes \texttt{SavedPostsPanel}, \texttt{SavedPostList}, \texttt{SavedPostsLink} labelled \textbf{Saved Posts}, \texttt{SavePostButton} with \textbf{Save}/\textbf{Saved} state, and \texttt{EmptySavedPostsState} with "No saved posts". \texttt{list\_post.md} also proposes \texttt{/forum/saved}; exact routes are not specified in proposals.
\item \textbf{Confirmed endpoint:} \texttt{POST /forum/posts/:post\_id/save} returns \texttt{SUCCESS 200} with \texttt{\{saved: boolean\}}; failures are \texttt{404} with \texttt{POST\_NOT\_FOUND}, \texttt{UNAUTHORIZED 401}, and \texttt{INTERNAL\_SERVER\_ERROR 500}.
\item \textbf{Database contract:} \texttt{studitopia.save\_forum\_post(puser\_id: uuid, ppost\_id: uuid)} returns \texttt{\{saved: boolean\}} and lists \texttt{UNKNOWN\_ERROR: 45000}, \texttt{POST\_NOT\_FOUND: 48001}, and \texttt{FORBIDDEN: 48010}.
\item \textbf{Detail state:} \texttt{GET /forum/posts/:post\_id} returns \texttt{saved: boolean} with \texttt{post}, \texttt{comments}, \texttt{true\_response}, and \texttt{user\_vote}; \texttt{404 POST\_NOT\_FOUND}, \texttt{401}, and \texttt{500} are defined.
\item \textbf{Gap:} \texttt{API.md} has no saved-list endpoint and no dedicated My Posts endpoint. \texttt{post\_management.md} marks the saved-list API as incomplete.
\item \textbf{Data fields:} \texttt{saved}, \texttt{post\_id}, \texttt{title}, \texttt{post\_type}, \texttt{author}, \texttt{vote\_score}, \texttt{comment\_count}, \texttt{solved}, and \texttt{created\_at} are defined by \texttt{API.md} list/detail contracts where applicable.
\item \textbf{Notifications and privacy:} No save-specific notification requirement is stated. Saved-post display remains subject to forum visibility, privacy, and moderation visibility rules.
\end{itemize}
\par\medskip

\hypertarget{sec:b11-source-alignment-notes}{}
\addcontentsline{toc}{subsubsection}{B11. Source Alignment Notes}
\subsubsection*{B11. Source Alignment Notes}
\par\medskip

\begin{itemize}
\item UC-03 covers browsing, searching, opening posts, and listing save as a permitted detail action; UC-19 extracts only the save, unsave, saved-list, and persistence behavior without changing UC-03.
\item UC-05 covers comment/answer creation and post-detail context where save may appear; it does not define saved-list behavior.
\item \texttt{FR-17.4} and \texttt{AC-17.4} allow save to appear among permission-based actions, but UC-19's requirements are the specific \texttt{FR-21.1}--\texttt{FR-21.4} and \texttt{AC-21.1}--\texttt{AC-21.4}.
\item \texttt{post\_management.md} recommends UI components and routes for saved posts but explicitly states that the saved-list API is incomplete in \texttt{API.md}; no endpoint is invented here.
\item No distinct include, extend, actor-specialization, or saved-list use-case ID is named for US-21, so the diagram contains only the selected use case and its primary actor.
\end{itemize}
\par\medskip

\pagebreak
\hypertarget{sec:uc-20-report-and-moderate-forum-content}{}
\addcontentsline{toc}{section}{UC-20 --- Report and Moderate Forum Content}
\section*{UC-20 --- Report and Moderate Forum Content}
\par\medskip

\textbf{Version:} \texttt{1.0}
\par\medskip

\textbf{Performed by:} [Nguyen Le Quoc Huy] | \textbf{Edited by:} [Nguyen Le Quoc Huy]
\par\medskip

\begingroup
\small
\renewcommand{\baselinestretch}{1.0}\selectfont
\renewcommand{\arraystretch}{1.15}
\setlength{\tabcolsep}{6pt}
\begin{longtable}{|>{\raggedright\arraybackslash}p{0.09\textwidth}|>{\raggedright\arraybackslash}p{0.84\textwidth}|}
\hline
\textbf{Field} & \textbf{Value} \\ \hline
\endhead
Use Case ID & UC-20 \\ \hline
Name & Report and Moderate Forum Content \\ \hline
Primary Actor & University Student or Moderator \\ \hline
Supporting Actors & Administrator \\ \hline
Feature Group & Forum Interaction \\ \hline
Priority & Not specified in proposals \\ \hline
Related User Stories & US-22 \\ \hline
Related Requirements & FR-22.1, FR-22.2, FR-22.3, FR-22.4, FR-22.5, FR-22.6; AC-22.1, AC-22.2, AC-22.3, AC-22.4, AC-22.5; NFR-SEC-1, NFR-SEC-6, NFR-SEC-9, NFR-SEC-10, NFR-USE-6, NFR-PRIV-7 \\ \hline
Source Documents & \texttt{proposals/userstory.md} §6.5; \texttt{proposals/functional\_requirement.md} US-22; \texttt{proposals/moderation\_queue.md}; \texttt{proposals/post\_management.md}; \texttt{proposals/list\_post.md}; \texttt{proposals/API.md}; \texttt{proposals/response\_code.md}; \texttt{proposals/non\_functional\_requirement.md}; \texttt{proposals/settings.md}; \texttt{proposals/feature\_flow.md}; \texttt{proposals/screen\_flow.md} \\ \hline
\end{longtable}
\endgroup
\par\medskip

\hypertarget{sec:part-a-use-case-model}{}
\addcontentsline{toc}{subsection}{Part A --- Use Case Model}
\subsection*{Part A --- Use Case Model}
\par\medskip

\hypertarget{sec:a1-focused-use-case-diagram}{}
\addcontentsline{toc}{subsubsection}{A1. Focused Use Case Diagram}
\subsubsection*{A1. Focused Use Case Diagram}
\par\medskip

\begin{center}
\includegraphics[width=0.85\textwidth,height=0.75\textheight,keepaspectratio]{"figures/UC-20-report-moderate-forum-content_1.png"}
\end{center}
\par\medskip

\textbf{Diagram notation.} Person shape = human actor; cloud = supporting system; blue = selected use case; solid line = association; labelled solid arrow = generalization; dashed arrow = include or extend.
\par\medskip

\hypertarget{sec:a2-actor-register}{}
\addcontentsline{toc}{subsubsection}{A2. Actor Register}
\subsubsection*{A2. Actor Register}
\par\medskip

\begingroup
\small
\renewcommand{\baselinestretch}{1.0}\selectfont
\renewcommand{\arraystretch}{1.15}
\setlength{\tabcolsep}{6pt}
\begin{longtable}{|>{\raggedright\arraybackslash}p{0.12\textwidth}|>{\raggedright\arraybackslash}p{0.15\textwidth}|>{\raggedright\arraybackslash}p{0.65\textwidth}|}
\hline
\textbf{Actor} & \textbf{Type} & \textbf{Description} \\ \hline
\endhead
University Student & Primary human actor & Reports a post, comment, or answer under the report categories in FR-22.2. \\ \hline
Moderator & Primary human actor & Reviews routed reports and applies permitted hide, delete, restore, or ban actions under FR-22.4. \\ \hline
Administrator & Supporting human actor & Is also authorized for moderation actions under NFR-SEC-10. \\ \hline
\end{longtable}
\endgroup
\par\medskip

\hypertarget{sec:a3-relationships}{}
\addcontentsline{toc}{subsubsection}{A3. Relationships}
\subsubsection*{A3. Relationships}
\par\medskip

\hypertarget{sec:a3-1-associations}{}
\paragraph*{A3.1 Associations}
\par\medskip

\begingroup
\small
\renewcommand{\baselinestretch}{1.0}\selectfont
\renewcommand{\arraystretch}{1.15}
\setlength{\tabcolsep}{6pt}
\begin{longtable}{|>{\raggedright\arraybackslash}p{0.13\textwidth}|>{\raggedright\arraybackslash}p{0.28\textwidth}|>{\raggedright\arraybackslash}p{0.10\textwidth}|>{\raggedright\arraybackslash}p{0.40\textwidth}|}
\hline
\textbf{Actor} & \textbf{Use Case} & \textbf{Type} & \textbf{Description} \\ \hline
\endhead
University Student & UC-20 Report and Moderate Forum Content & Primary & Submits a report on forum content. \\ \hline
Moderator & UC-20 Report and Moderate Forum Content & Primary & Reviews the moderation queue and applies policy actions. \\ \hline
Administrator & UC-20 Report and Moderate Forum Content & Supporting & Performs authorized moderation actions under NFR-SEC-10. \\ \hline
\end{longtable}
\endgroup
\par\medskip

\hypertarget{sec:a3-2-generalization-and-specialization}{}
\paragraph*{A3.2 Generalization and Specialization}
\par\medskip

\begingroup
\small
\renewcommand{\baselinestretch}{1.0}\selectfont
\renewcommand{\arraystretch}{1.15}
\setlength{\tabcolsep}{6pt}
\begin{longtable}{|>{\raggedright\arraybackslash}p{0.09\textwidth}|>{\raggedright\arraybackslash}p{0.09\textwidth}|>{\raggedright\arraybackslash}p{0.09\textwidth}|>{\raggedright\arraybackslash}p{0.65\textwidth}|}
\hline
\textbf{Child} & \textbf{Parent} & \textbf{Type} & \textbf{Inherited Behavior or Association} \\ \hline
\endhead
None & None & None & No actor generalization or specialization is named in the proposal sources for this story. \\ \hline
\end{longtable}
\endgroup
\par\medskip

\hypertarget{sec:a3-3-include}{}
\paragraph*{A3.3 Include}
\par\medskip

\begingroup
\small
\renewcommand{\baselinestretch}{1.0}\selectfont
\renewcommand{\arraystretch}{1.15}
\setlength{\tabcolsep}{6pt}
\begin{longtable}{|>{\raggedright\arraybackslash}p{0.10\textwidth}|>{\raggedright\arraybackslash}p{0.10\textwidth}|>{\raggedright\arraybackslash}p{0.72\textwidth}|}
\hline
\textbf{Base Use Case} & \textbf{Included Use Case} & \textbf{Meaning} \\ \hline
\endhead
None & None & No distinct included use case is named; \texttt{moderation\_queue.md} explicitly notes that there is no dedicated moderation use case. \\ \hline
\end{longtable}
\endgroup
\par\medskip

\hypertarget{sec:a3-4-extend}{}
\paragraph*{A3.4 Extend}
\par\medskip

\begingroup
\small
\renewcommand{\baselinestretch}{1.0}\selectfont
\renewcommand{\arraystretch}{1.15}
\setlength{\tabcolsep}{6pt}
\begin{longtable}{|>{\raggedright\arraybackslash}p{0.14\textwidth}|>{\raggedright\arraybackslash}p{0.10\textwidth}|>{\raggedright\arraybackslash}p{0.10\textwidth}|>{\raggedright\arraybackslash}p{0.58\textwidth}|}
\hline
\textbf{Extending Use Case} & \textbf{Base Use Case} & \textbf{Condition} & \textbf{Added Behavior} \\ \hline
\endhead
None & None & None & No distinct extending use case is named in the proposal sources for this story. \\ \hline
\end{longtable}
\endgroup
\par\medskip

\hypertarget{sec:a4-model-notes}{}
\addcontentsline{toc}{subsubsection}{A4. Model Notes}
\subsubsection*{A4. Model Notes}
\par\medskip

\begin{enumerate}
\item In scope: reporting a post, comment, or answer; routing reports with context; moderator/admin review; hide, delete, restore, or ban actions; action records; notifications; and restriction enforcement.
\item Out of scope: creating posts and comments, profile/message reporting under US-40, and modifying UC-03 or UC-05.
\item The four exact report categories are spam, harassment, inappropriate content, and academic dishonesty.
\item \texttt{moderation\_queue.md} states that \texttt{API.md} defines only \texttt{POST /forum/reports}; no queue, action, history, or status endpoints are defined.
\item No new I-xx, E-xx, or G-xx identifier is created; the queue and moderation record are UI/data recommendations, not named use cases in the proposal sources.
\end{enumerate}
\par\medskip

\hypertarget{sec:part-b-use-case-specification}{}
\addcontentsline{toc}{subsection}{Part B --- Use Case Specification}
\subsection*{Part B --- Use Case Specification}
\par\medskip

\hypertarget{sec:b1-brief-description}{}
\addcontentsline{toc}{subsubsection}{B1. Brief Description}
\subsubsection*{B1. Brief Description}
\par\medskip

A university student reports inappropriate forum content, and an authorized moderator or administrator reviews the report and applies policy actions so that the forum remains safe and academically useful. Reports carry target context and are routed to a moderation queue.
\par\medskip

The trigger is selecting \textbf{Report} on a post, comment, or answer, or opening the moderation queue as an authorized reviewer. The successful outcome is a recorded report with \texttt{report\_id}, a moderation decision recorded under FR-22.5, notification to the affected user when appropriate, and enforced restrictions where applicable.
\par\medskip

\hypertarget{sec:b2-preconditions}{}
\addcontentsline{toc}{subsubsection}{B2. Preconditions}
\subsubsection*{B2. Preconditions}
\par\medskip

\begin{enumerate}
\item The reporting student or reviewing moderator/administrator is authenticated.
\item The reported post, comment, or answer exists and is visible to the reporter.
\item The reporter selects one of spam, harassment, inappropriate content, or academic dishonesty.
\item The reviewer is an authorized moderator or administrator under NFR-SEC-10.
\end{enumerate}
\par\medskip

\hypertarget{sec:b3-postconditions}{}
\addcontentsline{toc}{subsubsection}{B3. Postconditions}
\subsubsection*{B3. Postconditions}
\par\medskip

\begin{enumerate}
\item \textbf{Success:} A report is recorded with target context and returned as \texttt{report\_id}, then routed to the moderation queue.
\item \textbf{Success:} An authorized reviewer applies a permitted action and the action is recorded with notification to the affected user when appropriate.
\item \textbf{Success:} Moderation restrictions cannot be bypassed through new posts or comments.
\item \textbf{Failure:} No invalid or unauthorized moderation change is applied, and the actor receives an error.
\end{enumerate}
\par\medskip

\hypertarget{sec:b4-main-success-scenario-basic-flow}{}
\addcontentsline{toc}{subsubsection}{B4. Main Success Scenario (Basic Flow)}
\subsubsection*{B4. Main Success Scenario (Basic Flow)}
\par\medskip

\begin{enumerate}
\item The student opens a post detail view containing a post, comment, or answer to report.
\item The system shows \textbf{Report} among the permitted actions.
\item The student selects \textbf{Report} and the system opens the report dialog with target context.
\item The student selects a category: spam, harassment, inappropriate content, or academic dishonesty.
\item The student optionally enters a description and submits the report.
\item The backend validates the target, category, and reporter context and records the report.
\item The backend returns \texttt{201} with \texttt{\{report\_id\}} and routes the report to the moderation queue with contextual information.
\item An authorized moderator or administrator reviews the report context.
\item The reviewer selects a policy action: hide, delete, restore, or ban content or users.
\item The system records the moderation action and notifies the affected user when appropriate.
\item The system enforces any resulting restriction on new posts and comments.
\end{enumerate}
\par\medskip

\hypertarget{sec:b5-alternative-flows}{}
\addcontentsline{toc}{subsubsection}{B5. Alternative Flows}
\subsubsection*{B5. Alternative Flows}
\par\medskip

\hypertarget{sec:a1-report-a-comment-or-answer}{}
\paragraph*{A1. Report a comment or answer}
\par\medskip

At basic flow steps 1-6, the reporter targets a comment or answer. The request uses \texttt{target\_type: 'comment'} and the corresponding \texttt{target\_id}; a missing target produces \texttt{COMMENT\_NOT\_FOUND} (\texttt{48002}).
\par\medskip

\hypertarget{sec:a2-review-without-description}{}
\paragraph*{A2. Review without description}
\par\medskip

At basic flow step 5, the reporter leaves the optional description empty. The backend accepts the report when \texttt{target\_type}, \texttt{target\_id}, and a valid \texttt{category} are present.
\par\medskip

\hypertarget{sec:a3-restore-moderated-content}{}
\paragraph*{A3. Restore moderated content}
\par\medskip

At basic flow step 9, the reviewer restores previously hidden or deleted content where supported by forum policy.
\par\medskip

\hypertarget{sec:a4-notify-affected-user}{}
\paragraph*{A4. Notify affected user}
\par\medskip

At basic flow step 10, the moderation record indicates that an affected user should be notified. The system notifies that user when appropriate under AC-22.4.
\par\medskip

\hypertarget{sec:a5-restriction-enforcement}{}
\paragraph*{A5. Restriction enforcement}
\par\medskip

At basic flow step 11, a moderated user attempts a new post or comment. The backend blocks the bypass attempt and the UI shows a restriction notice or action error.
\par\medskip

\hypertarget{sec:b6-exception-flows}{}
\addcontentsline{toc}{subsubsection}{B6. Exception Flows}
\subsubsection*{B6. Exception Flows}
\par\medskip

\hypertarget{sec:e1-invalid-report-category}{}
\paragraph*{E1. Invalid report category}
\par\medskip

The backend rejects an unsupported category with \texttt{400 INVALID\_REPORT\_CATEGORY} (\texttt{48009}) and the report is not recorded.
\par\medskip

\hypertarget{sec:e2-target-is-unavailable}{}
\paragraph*{E2. Target is unavailable}
\par\medskip

A missing post or comment produces \texttt{404 POST\_NOT\_FOUND} (\texttt{48001}) or \texttt{COMMENT\_NOT\_FOUND} (\texttt{48002}), and the report is not recorded.
\par\medskip

\hypertarget{sec:e3-unauthenticated-report}{}
\paragraph*{E3. Unauthenticated report}
\par\medskip

An unauthenticated report request produces \texttt{UNAUTHORIZED} under \texttt{401}; the system requires login before submitting.
\par\medskip

\hypertarget{sec:e4-unauthorized-moderation-access}{}
\paragraph*{E4. Unauthorized moderation access}
\par\medskip

A user who is not an authorized moderator or administrator is denied moderation access under NFR-SEC-10 and server-side checks under NFR-SEC-6. \texttt{moderation\_queue.md} proposes a permission-denied state; no moderation-action HTTP contract is defined in \texttt{API.md}.
\par\medskip

\hypertarget{sec:e5-duplicate-report}{}
\paragraph*{E5. Duplicate report}
\par\medskip

\texttt{API.md} requires preventing duplicate reports from the same user for the same target within a short window, but defines no error code for that case. The UI must not claim a specific duplicate code.
\par\medskip

\hypertarget{sec:e6-server-or-database-failure}{}
\paragraph*{E6. Server or database failure}
\par\medskip

The backend reports \texttt{INTERNAL\_SERVER\_ERROR} under \texttt{500} with \texttt{(DATABASE\_ERROR | UNKNOWN\_ERROR)[]}; the report or moderation action is not completed.
\par\medskip

\hypertarget{sec:b7-business-rules}{}
\addcontentsline{toc}{subsubsection}{B7. Business Rules}
\subsubsection*{B7. Business Rules}
\par\medskip

\begin{enumerate}
\item Users can report a post, comment, or answer under FR-22.1 and AC-22.1.
\item The exact report categories are spam, harassment, inappropriate content, and academic dishonesty under FR-22.2 and AC-22.1.
\item Reports are routed to a moderation queue with contextual information under FR-22.3 and AC-22.2.
\item Moderators can hide, delete, restore, or ban content and users according to forum policy under FR-22.4 and AC-22.3.
\item Moderation actions are recorded and affected users are notified when appropriate under FR-22.5 and AC-22.4.
\item Moderated users cannot bypass restrictions through new posts or comments under FR-22.6 and AC-22.5.
\item Moderation actions are restricted to authorized moderators and administrators under NFR-SEC-10.
\item Reporting and related actions are rate-limited under NFR-SEC-9; no numeric limit is specified in proposals.
\end{enumerate}
\par\medskip

\hypertarget{sec:b8-special-requirements}{}
\addcontentsline{toc}{subsubsection}{B8. Special Requirements}
\subsubsection*{B8. Special Requirements}
\par\medskip

\begingroup
\small
\renewcommand{\baselinestretch}{1.0}\selectfont
\renewcommand{\arraystretch}{1.15}
\setlength{\tabcolsep}{6pt}
\begin{longtable}{|>{\raggedright\arraybackslash}p{0.12\textwidth}|>{\raggedright\arraybackslash}p{0.81\textwidth}|}
\hline
\textbf{ID} & \textbf{Requirement} \\ \hline
\endhead
NFR-SEC-1 & Authentication is required before accessing protected report and moderation surfaces. \\ \hline
NFR-SEC-6 & Server-side authorization checks must prevent unauthorized moderation access. \\ \hline
NFR-SEC-9 & Reporting, posting, commenting, and voting must be rate-limited to reduce spam and abuse. \\ \hline
NFR-SEC-10 & Moderation actions must be restricted to authorized moderators and administrators. \\ \hline
NFR-USE-6 & Reporting and moderation actions must be reachable within at most 2 interactions from a post or comment. \\ \hline
NFR-PRIV-7 & Privacy settings must apply consistently to forum posts, comments, member lists, and search results. \\ \hline
Design constraint & Moderation actions must produce a moderation record; no queue, action, or history endpoint is defined in \texttt{API.md}. \\ \hline
\end{longtable}
\endgroup
\par\medskip

\hypertarget{sec:b9-extension-points}{}
\addcontentsline{toc}{subsubsection}{B9. Extension Points}
\subsubsection*{B9. Extension Points}
\par\medskip

\begingroup
\small
\renewcommand{\baselinestretch}{1.0}\selectfont
\renewcommand{\arraystretch}{1.15}
\setlength{\tabcolsep}{6pt}
\begin{longtable}{|>{\raggedright\arraybackslash}p{0.39\textwidth}|>{\raggedright\arraybackslash}p{0.54\textwidth}|}
\hline
\textbf{Extension Point} & \textbf{Location in Flow of Events} \\ \hline
\endhead
Report target selection & Basic flow steps 1-4; alternative A1 \\ \hline
Category and description entry & Basic flow steps 4-5; exceptions E1 and E5 \\ \hline
Queue review & Basic flow step 8; exception E4 \\ \hline
Moderation action & Basic flow step 9; alternatives A3-A4 \\ \hline
Restriction enforcement & Basic flow step 11; alternative A5 \\ \hline
Failure recovery & Exceptions E1-E6 \\ \hline
\end{longtable}
\endgroup
\par\medskip

\hypertarget{sec:b10-interface-and-data-notes}{}
\addcontentsline{toc}{subsubsection}{B10. Interface and Data Notes}
\subsubsection*{B10. Interface and Data Notes}
\par\medskip

\begin{itemize}
\item \textbf{Confirmed flow:} \texttt{feature\_flow.md} shows \texttt{Post → Report / Moderate}; \texttt{moderation\_queue.md} shows \texttt{Report Content → Moderation Queue → Moderator Action → Record / Notify / Enforce Restriction}.
\item \textbf{Confirmed categories:} Report categories are spam, harassment, inappropriate content, and academic dishonesty under FR-22.2 and AC-22.1.
\item \textbf{Recommendation UI:} \texttt{ReportPostButton} is labelled \textbf{Report}; \texttt{ReportContentDialog} / \texttt{ReportPostDialog} provides target context, category, optional description, and submit/cancel; \texttt{ReportSubmitButton} is labelled \textbf{Submit report}. These are source recommendations, not FR/AC fields.
\item \textbf{Confirmed endpoint:} \texttt{POST /forum/reports} uses \texttt{\{target\_type: 'post' | 'comment', target\_id: string, category: string, description?: string\}} and returns \texttt{201 \{report\_id\}}.
\item \textbf{Error codes:} \texttt{400 INVALID\_REPORT\_CATEGORY} (\texttt{48009}), \texttt{404 POST\_NOT\_FOUND} (\texttt{48001}) or \texttt{COMMENT\_NOT\_FOUND} (\texttt{48002}), \texttt{401}, and \texttt{500}. \texttt{studitopia.report\_forum\_content} also lists \texttt{FORBIDDEN} (\texttt{48010}) and duplicate-report prevention with no specified error code.
\item \textbf{Gap:} No queue, action, history, or status endpoints are defined in \texttt{API.md}. Hide, delete, restore, and ban actions have requirement coverage in FR-22.4 but no API contract.
\item \textbf{Authorization:} Only authorized moderators/admins may act, per NFR-SEC-10, with server-side checks under NFR-SEC-6.
\item \textbf{Data fields:} Confirmed fields are \texttt{report\_id}, \texttt{target\_type}, \texttt{target\_id}, \texttt{category}, and optional \texttt{description}. \texttt{reporter\_id} appears in the database function; \texttt{created\_at}, \texttt{status}, \texttt{moderation\_action}, \texttt{action\_reason}, and \texttt{affected\_user\_notified} are recommendation/gap fields in \texttt{moderation\_queue.md}.
\item \textbf{Moderation record:} FR-22.5 and AC-22.4 require recorded actions and notification to affected users when appropriate; no history API is defined.
\item \textbf{Privacy:} NFR-PRIV-7 applies consistently to posts, comments, member lists, and search results.
\end{itemize}
\par\medskip

\hypertarget{sec:b11-source-alignment-notes}{}
\addcontentsline{toc}{subsubsection}{B11. Source Alignment Notes}
\subsubsection*{B11. Source Alignment Notes}
\par\medskip

\begin{itemize}
\item UC-03 lists \textbf{report} as a permitted action in the post-detail flow but does not define report submission, queue review, moderation actions, records, or restriction enforcement.
\item UC-05 mentions the same \texttt{POST /forum/reports} contract and the moderation API gap as interface context; it does not cover reporting or moderation behavior.
\item \texttt{moderation\_queue.md} explicitly states that there is no dedicated moderation use case in the cited use-case files, so UC-20 introduces no new I-xx/E-xx/G-xx relationship.
\item \texttt{FR-40.1}--\texttt{FR-40.5} and \texttt{AC-40.1}--\texttt{AC-40.4} cover broader inappropriate-behavior reporting, but this case is scoped to US-22 forum content.
\item \texttt{API.md} supports only \texttt{target\_type: 'post' | 'comment'}; report surfaces beyond posts and comments lack report API support.
\item No queue, action, history, or status endpoint is defined in \texttt{API.md}, and no report-status enum is specified; do not infer one from UI recommendations.
\end{itemize}
\par\medskip

\par\medskip
\hypertarget{sec:ui-prototype-specimens}{}
\addcontentsline{toc}{subsection}{UI Prototype Specimens}
\subsection*{UI Prototype Specimens}
\par\medskip
\begin{figure}[H]
\centering
\includegraphics[width=0.75\textwidth,height=0.42\textheight,keepaspectratio]{"../../prototype/profile_p5_report_a_user.png"}
\par\vspace{2pt}
{\footnotesize \textbf{UI Prototype:} Profile: P5 Report a user}
\end{figure}
\par\smallskip

\end{document}
